% Recovered from the previously completed PDF; no new mathematical content.
% Prose is editable; formula boxes retain original symbol positions.
% Build from repository root: xelatex -no-shell-escape -halt-on-error -output-directory=... print/fudan/NAME.tex
\documentclass[11pt,a4paper]{article}
\usepackage[margin=23mm,headheight=14pt,headsep=8.5mm,footskip=10mm]{geometry}
\usepackage{amsmath,amssymb}
\DeclareMathAlphabet{\mathsf}{OT1}{cmss}{m}{n}
\usepackage[no-math]{fontspec}
\usepackage{xeCJK}
\setmainfont{Latin Modern Roman}
\setCJKmainfont{FandolSong-Regular.otf}[Path=assets/fonts/fandol/,BoldFont=FandolSong-Bold.otf,ItalicFont=FandolSong-Regular.otf,Script=CustomDefault]
\xeCJKsetup{PunctStyle=quanjiao,CheckSingle=true}
\usepackage{setspace}\setstretch{1.12}
\setlength{\parindent}{2em}\setlength{\parskip}{3pt}
\emergencystretch=6em
\hfuzz=9pt
\usepackage{titlesec}
\titleformat{\section}{\centering\normalfont\Large\bfseries}{}{0pt}{}
\titlespacing*{\section}{0pt}{1.8\baselineskip}{1.0\baselineskip}
\usepackage{fancyhdr}\pagestyle{fancy}\fancyhf{}
\fancyhead[L]{\small 数学手记·高等代数}\fancyhead[R]{\small 15A08–15A12、16S01–16S02}
\fancyfoot[C]{\thepage}\renewcommand{\headrulewidth}{0.4pt}
\usepackage[hidelinks,unicode]{hyperref}
\hypersetup{pdftitle={15A08–15A12、16S01–16S02：题目、解答与小结},pdfauthor={Vincent Arqain}}
\newcommand{\RB}[4]{\raisebox{0pt}[#2bp][#3bp]{\hbox to #1bp{#4\hss}}}
\newcommand{\RG}[5]{\hbox to 0pt{\kern#1bp\raisebox{#2bp}[0pt][0pt]{\fontsize{#3bp}{#3bp}\csname font#4\endcsname\char#5}\hss}}
\newcommand{\RT}[4]{\hbox to 0pt{\kern#1bp\raisebox{#2bp}[0pt][0pt]{\fontsize{#3bp}{#3bp}\selectfont#4}\hss}}
\newcommand{\RS}[4]{\hbox to 0pt{\kern#1bp\raisebox{#2bp}[0pt][0pt]{\fontsize{#3bp}{#3bp}\selectfont\ensuremath{#4}}\hss}}
\newcommand{\RR}[4]{\hbox to 0pt{\kern#1bp\raisebox{#2bp}[0pt][0pt]{\rule{#3bp}{#4bp}}\hss}}
\newcommand{\fontmi}{\usefont{OML}{cmm}{m}{it}}
\newcommand{\fontsy}{\usefont{OMS}{cmsy}{m}{n}}
\newcommand{\fontex}{\usefont{OMX}{cmex}{m}{n}}
\newcommand{\fontrm}{\usefont{OT1}{cmr}{m}{n}}
\newcommand{\fontbf}{\usefont{OT1}{cmr}{bx}{n}}
\newcommand{\fontam}{\usefont{U}{msa}{m}{n}}
\newcommand{\fontbb}{\usefont{U}{msb}{m}{n}}
\widowpenalty=10000\clubpenalty=10000\raggedbottom
\begin{document}
\begin{center}{\LARGE\bfseries 15A08–15A12、16S01–16S02\par}\vspace{8pt}复旦大学高等代数习题集·题目与解答\end{center}
% Original page 1.
复旦大学高等代数习题集题目与证明

题目选自谢启鸿《复旦大学高等代数习题集》（2023版），§1.4与§1.5，第20—21、24页。证明与小结为本次整理。

记号：\allowbreak{}\hskip.1em \RB{20.755}{7.121}{2.585}{%
\RG{0.000}{-0.004}{7.9701}{rm}{107}\RG{4.782}{-0.004}{7.9701}{rm}{101}\RG{9.020}{-0.004}{7.9701}{rm}{114}\RG{14.346}{-0.004}{7.9701}{mi}{39}%
}\hskip.1em\allowbreak{} 、\allowbreak{}\hskip.1em \RB{19.471}{7.016}{2.585}{%
\RG{0.000}{-0.004}{7.9701}{rm}{73}\RG{3.446}{-0.004}{7.9701}{rm}{109}\RG{13.063}{-0.004}{7.9701}{mi}{39}%
}\hskip.1em\allowbreak{} 分别表示核与像，\allowbreak{}\hskip.1em \RB{68.010}{7.655}{2.890}{%
\RG{0.000}{-0.004}{7.9701}{mi}{114}\RG{4.572}{-0.004}{7.9701}{rm}{40}\RG{8.284}{-0.004}{7.9701}{mi}{39}\RG{14.528}{-0.004}{7.9701}{rm}{41}\RG{20.885}{-0.004}{7.9701}{rm}{61}\RG{30.965}{-0.004}{7.9701}{rm}{100}\RG{36.273}{-0.004}{7.9701}{rm}{105}\RG{38.926}{-0.004}{7.9701}{rm}{109}\RG{48.539}{-0.004}{7.9701}{rm}{73}\RG{51.985}{-0.004}{7.9701}{rm}{109}\RG{61.602}{-0.004}{7.9701}{mi}{39}%
}\hskip.1em\allowbreak{} ；\allowbreak{}\hskip.1em \RB{8.868}{7.016}{1.999}{%
\RG{0.000}{-0.004}{7.9701}{mi}{73}\RG{4.196}{-1.435}{7.9701}{mi}{110}%
}\hskip.1em\allowbreak{} 表示\allowbreak{}\hskip.1em \RB{5.950}{4.715}{0.609}{%
\RG{0.000}{-0.004}{7.9701}{mi}{110}%
}\hskip.1em\allowbreak{} 阶单位矩阵。

\section*{2015A08}
设\allowbreak{}\hskip.1em \RB{7.122}{5.307}{2.860}{%
\RG{0.000}{-0.004}{10.9091}{mi}{39}%
}\hskip.1em\allowbreak{} 是\allowbreak{}\hskip.1em \RB{6.555}{5.307}{0.613}{%
\RG{0.000}{-0.004}{10.9091}{mi}{110}%
}\hskip.1em\allowbreak{} 维线性空间\allowbreak{}\hskip.1em \RB{8.715}{7.947}{0.733}{%
\RG{0.000}{-0.004}{10.9091}{mi}{86}%
}\hskip.1em\allowbreak{} 上的线性变换，满足\allowbreak{}\hskip.1em \RB{34.012}{7.751}{2.860}{%
\RG{0.000}{-0.004}{10.9091}{mi}{39}\RG{6.997}{3.955}{7.9701}{mi}{109}\RG{17.290}{-0.004}{10.9091}{rm}{61}\RG{28.625}{-0.004}{10.9091}{rm}{48}%
}\hskip.1em\allowbreak{} ，其中\allowbreak{}\hskip.1em \RB{19.438}{5.307}{2.620}{%
\RG{0.000}{-0.004}{10.9091}{mi}{109}\RG{9.349}{-0.004}{10.9091}{mi}{59}\RG{14.127}{-0.004}{10.9091}{mi}{113}%
}\hskip.1em\allowbreak{} 为正整数，\allowbreak{}\hskip.1em \RB{50.205}{7.751}{2.620}{%
\RG{0.000}{-0.004}{10.9091}{mi}{110}\RG{9.406}{-0.004}{10.9091}{rm}{61}\RG{20.745}{-0.004}{10.9091}{mi}{109}\RG{30.095}{-0.004}{10.9091}{mi}{113}\RG{36.068}{-0.004}{10.9091}{rm}{43}\RG{45.232}{-0.004}{10.9091}{rm}{49}%
}\hskip.1em\allowbreak{} 。

证明：

\[
\RB{101.861}{8.067}{2.860}{%
\RG{0.000}{-0.004}{10.9091}{rm}{100}\RG{5.935}{-0.004}{10.9091}{rm}{105}\RG{8.902}{-0.004}{10.9091}{rm}{109}\RG{19.669}{-0.004}{10.9091}{rm}{73}\RG{23.520}{-0.004}{10.9091}{rm}{109}\RG{34.300}{-0.004}{10.9091}{mi}{39}\RG{44.327}{-0.004}{10.9091}{sy}{20}\RG{55.843}{-0.004}{10.9091}{mi}{110}\RG{64.643}{-0.004}{10.9091}{sy}{0}\RG{75.552}{-0.004}{10.9091}{mi}{113}\RG{83.104}{-0.004}{10.9091}{sy}{0}\RG{94.013}{-0.004}{10.9091}{rm}{49}\RG{99.354}{-0.004}{10.9091}{mi}{58}%
}
\]
证明：记\allowbreak{}\hskip.1em \RB{53.991}{9.255}{3.615}{%
\RG{0.000}{-0.001}{10.9091}{mi}{75}\RG{9.087}{-1.637}{7.9701}{mi}{106}\RG{16.102}{-0.001}{10.9091}{rm}{61}\RG{27.436}{-0.001}{10.9091}{rm}{107}\RG{32.793}{-0.001}{10.9091}{rm}{101}\RG{37.538}{-0.001}{10.9091}{rm}{114}\RG{43.500}{-0.001}{10.9091}{mi}{39}\RG{50.498}{3.958}{7.9701}{mi}{106}%
}\hskip.1em\allowbreak{} （\allowbreak{}\hskip.1em \RB{48.986}{7.754}{2.726}{%
\RG{0.000}{-0.001}{10.9091}{rm}{48}\RG{8.371}{-0.001}{10.9091}{sy}{20}\RG{19.886}{-0.001}{10.9091}{mi}{106}\RG{27.949}{-0.001}{10.9091}{sy}{20}\RG{39.465}{-0.001}{10.9091}{mi}{109}%
}\hskip.1em\allowbreak{} ），其中\allowbreak{}\hskip.1em \RB{43.713}{8.681}{3.228}{%
\RG{0.000}{-0.001}{10.9091}{mi}{75}\RG{9.088}{-1.637}{7.9701}{rm}{48}\RG{16.420}{-0.001}{10.9091}{rm}{61}\RG{27.759}{-0.001}{10.9091}{sy}{102}\RG{33.214}{-0.001}{10.9091}{rm}{48}\RG{38.555}{-0.001}{10.9091}{sy}{103}%
}\hskip.1em\allowbreak{} ，并令\allowbreak{}\hskip.1em \RB{62.775}{8.070}{2.857}{%
\RG{0.000}{-0.001}{10.9091}{mi}{100}\RG{8.581}{-0.001}{10.9091}{rm}{61}\RG{19.915}{-0.001}{10.9091}{rm}{100}\RG{25.850}{-0.001}{10.9091}{rm}{105}\RG{28.817}{-0.001}{10.9091}{rm}{109}\RG{39.595}{-0.001}{10.9091}{rm}{107}\RG{44.941}{-0.001}{10.9091}{rm}{101}\RG{49.686}{-0.001}{10.9091}{rm}{114}\RG{55.654}{-0.001}{10.9091}{mi}{39}%
}\hskip.1em\allowbreak{} 。对每个\allowbreak{}\hskip.1em \RB{48.986}{7.754}{2.726}{%
\RG{0.000}{-0.001}{10.9091}{rm}{49}\RG{8.371}{-0.001}{10.9091}{sy}{20}\RG{19.886}{-0.001}{10.9091}{mi}{106}\RG{27.949}{-0.001}{10.9091}{sy}{20}\RG{39.465}{-0.001}{10.9091}{mi}{109}%
}\hskip.1em\allowbreak{} ，有

\[
\RB{165.771}{8.681}{4.350}{%
\RG{0.000}{-0.001}{10.9091}{mi}{39}\RG{6.997}{-0.001}{10.9091}{rm}{40}\RG{11.152}{-0.001}{10.9091}{mi}{75}\RG{24.223}{-0.001}{10.9091}{rm}{41}\RG{31.407}{-0.001}{10.9091}{sy}{18}\RG{42.922}{-0.001}{10.9091}{mi}{75}\RG{65.831}{-0.001}{10.9091}{mi}{59}\RG{92.435}{-0.001}{10.9091}{rm}{107}\RG{97.780}{-0.001}{10.9091}{rm}{101}\RG{102.525}{-0.001}{10.9091}{rm}{114}\RG{106.682}{-0.001}{10.9091}{rm}{40}\RG{110.830}{-0.001}{10.9091}{mi}{39}\RG{117.827}{-0.001}{10.9091}{sy}{106}\RG{131.351}{-0.001}{10.9091}{rm}{41}\RG{138.529}{-0.001}{10.9091}{rm}{61}\RG{149.874}{-0.001}{10.9091}{mi}{75}\RG{163.264}{-0.001}{10.9091}{mi}{58}\RG{20.239}{-1.637}{7.9701}{mi}{106}\RG{52.010}{-1.637}{7.9701}{mi}{106}\RG{55.541}{-1.637}{7.9701}{sy}{0}\RG{61.528}{-1.637}{7.9701}{rm}{49}\RG{120.858}{-1.637}{7.9701}{mi}{75}\RG{158.961}{-1.637}{7.9701}{rm}{49}\RG{127.350}{-2.742}{5.9776}{mi}{106}%
}
\]
由秩与零度定理，

\[
\RB{192.214}{8.685}{3.612}{%
\RG{0.000}{0.003}{10.9091}{rm}{100}\RG{5.935}{0.003}{10.9091}{rm}{105}\RG{8.902}{0.003}{10.9091}{rm}{109}\RG{19.675}{0.003}{10.9091}{mi}{75}\RG{35.775}{0.003}{10.9091}{rm}{61}\RG{47.115}{0.003}{10.9091}{mi}{100}\RG{55.089}{0.003}{10.9091}{rm}{43}\RG{65.824}{0.003}{10.9091}{rm}{100}\RG{71.758}{0.003}{10.9091}{rm}{105}\RG{74.725}{0.003}{10.9091}{rm}{109}\RG{85.495}{0.003}{10.9091}{mi}{39}\RG{92.493}{0.003}{10.9091}{rm}{40}\RG{96.646}{0.003}{10.9091}{mi}{75}\RG{109.718}{0.003}{10.9091}{rm}{41}\RG{116.903}{0.003}{10.9091}{sy}{20}\RG{128.417}{0.003}{10.9091}{mi}{100}\RG{136.392}{0.003}{10.9091}{rm}{43}\RG{147.126}{0.003}{10.9091}{rm}{100}\RG{153.061}{0.003}{10.9091}{rm}{105}\RG{156.028}{0.003}{10.9091}{rm}{109}\RG{166.798}{0.003}{10.9091}{mi}{75}\RG{189.706}{0.003}{10.9091}{mi}{58}\RG{28.762}{-1.634}{7.9701}{mi}{106}\RG{105.735}{-1.634}{7.9701}{mi}{106}\RG{175.886}{-1.634}{7.9701}{mi}{106}\RG{179.417}{-1.634}{7.9701}{sy}{0}\RG{185.405}{-1.634}{7.9701}{rm}{49}%
}
\]
逐次应用此不等式，得到\allowbreak{}\hskip.1em \RB{58.340}{8.069}{3.616}{%
\RG{0.000}{-0.002}{10.9091}{rm}{100}\RG{5.935}{-0.002}{10.9091}{rm}{105}\RG{8.902}{-0.002}{10.9091}{rm}{109}\RG{19.674}{-0.002}{10.9091}{mi}{75}\RG{28.762}{-1.638}{7.9701}{mi}{106}\RG{35.775}{-0.002}{10.9091}{sy}{20}\RG{47.291}{-0.002}{10.9091}{mi}{106}\RG{52.320}{-0.002}{10.9091}{mi}{100}%
}\hskip.1em\allowbreak{} 。又因\allowbreak{}\hskip.1em \RB{34.012}{7.753}{2.858}{%
\RG{0.000}{-0.002}{10.9091}{mi}{39}\RG{6.997}{3.956}{7.9701}{mi}{109}\RG{17.290}{-0.002}{10.9091}{rm}{61}\RG{28.625}{-0.002}{10.9091}{rm}{48}%
}\hskip.1em\allowbreak{} ，故\allowbreak{}\hskip.1em \RB{39.434}{7.949}{2.211}{%
\RG{0.000}{-0.002}{10.9091}{mi}{75}\RG{9.088}{-1.638}{7.9701}{mi}{109}\RG{19.380}{-0.002}{10.9091}{rm}{61}\RG{30.719}{-0.002}{10.9091}{mi}{86}%
}\hskip.1em\allowbreak{} ，从而

\[
\RB{136.061}{8.074}{2.614}{%
\RG{0.000}{0.003}{10.9091}{mi}{109}\RG{9.349}{0.003}{10.9091}{mi}{113}\RG{16.895}{0.003}{10.9091}{rm}{43}\RG{27.630}{0.003}{10.9091}{rm}{49}\RG{35.997}{0.003}{10.9091}{rm}{61}\RG{47.337}{0.003}{10.9091}{mi}{110}\RG{56.744}{0.003}{10.9091}{rm}{61}\RG{68.078}{0.003}{10.9091}{rm}{100}\RG{74.013}{0.003}{10.9091}{rm}{105}\RG{76.980}{0.003}{10.9091}{rm}{109}\RG{87.756}{0.003}{10.9091}{mi}{75}\RG{107.136}{0.003}{10.9091}{sy}{20}\RG{118.652}{0.003}{10.9091}{mi}{109}\RG{128.001}{0.003}{10.9091}{mi}{100}\RG{133.554}{0.003}{10.9091}{mi}{58}\RG{96.844}{-1.634}{7.9701}{mi}{109}%
}
\]
于是\allowbreak{}\hskip.1em \RB{67.039}{9.611}{4.337}{%
\RG{0.000}{-0.003}{10.9091}{mi}{100}\RG{8.581}{-0.003}{10.9091}{sy}{21}\RG{20.096}{-0.003}{10.9091}{mi}{113}\RG{27.648}{-0.003}{10.9091}{rm}{43}\RG{39.467}{-3.765}{7.9701}{mi}{109}\RG{40.947}{4.293}{7.9701}{rm}{49}\RG{50.394}{-0.003}{10.9091}{mi}{62}\RG{61.728}{-0.003}{10.9091}{mi}{113}\RR{39.467}{2.725}{6.809}{0.362}%
}\hskip.1em\allowbreak{} 。由于\allowbreak{}\hskip.1em \RB{6.020}{8.068}{0.612}{%
\RG{0.000}{-0.003}{10.9091}{mi}{100}%
}\hskip.1em\allowbreak{} 为整数，必有\allowbreak{}\hskip.1em \RB{43.355}{8.068}{2.619}{%
\RG{0.000}{-0.003}{10.9091}{mi}{100}\RG{8.581}{-0.003}{10.9091}{sy}{21}\RG{20.095}{-0.003}{10.9091}{mi}{113}\RG{27.647}{-0.003}{10.9091}{rm}{43}\RG{38.382}{-0.003}{10.9091}{rm}{49}%
}\hskip.1em\allowbreak{} ，故

\[
\RB{141.489}{8.073}{2.855}{%
\RG{0.000}{0.002}{10.9091}{rm}{100}\RG{5.935}{0.002}{10.9091}{rm}{105}\RG{8.902}{0.002}{10.9091}{rm}{109}\RG{19.669}{0.002}{10.9091}{rm}{73}\RG{23.520}{0.002}{10.9091}{rm}{109}\RG{34.299}{0.002}{10.9091}{mi}{39}\RG{44.327}{0.002}{10.9091}{rm}{61}\RG{55.665}{0.002}{10.9091}{mi}{110}\RG{64.465}{0.002}{10.9091}{sy}{0}\RG{75.375}{0.002}{10.9091}{mi}{100}\RG{83.955}{0.002}{10.9091}{sy}{20}\RG{95.471}{0.002}{10.9091}{mi}{110}\RG{104.271}{0.002}{10.9091}{sy}{0}\RG{115.180}{0.002}{10.9091}{mi}{113}\RG{122.732}{0.002}{10.9091}{sy}{0}\RG{133.641}{0.002}{10.9091}{rm}{49}\RG{138.982}{0.002}{10.9091}{mi}{58}%
}
\]
小结：幂零性使核空间在有限次迭代后覆盖整个空间，而每次核维数的增加至多为\allowbreak{}\hskip.1em \RB{42.856}{8.074}{2.854}{%
\RG{0.000}{0.003}{10.9091}{rm}{100}\RG{5.935}{0.003}{10.9091}{rm}{105}\RG{8.902}{0.003}{10.9091}{rm}{109}\RG{19.669}{0.003}{10.9091}{rm}{107}\RG{25.025}{0.003}{10.9091}{rm}{101}\RG{29.771}{0.003}{10.9091}{rm}{114}\RG{35.735}{0.003}{10.9091}{mi}{39}%
}\hskip.1em\allowbreak{} 。一般地，\allowbreak{}\hskip.1em \RB{34.012}{7.757}{2.854}{%
\RG{0.000}{0.003}{10.9091}{mi}{39}\RG{6.997}{3.962}{7.9701}{mi}{109}\RG{17.290}{0.003}{10.9091}{rm}{61}\RG{28.625}{0.003}{10.9091}{rm}{48}%
}\hskip.1em\allowbreak{} 蕴含\allowbreak{}\hskip.1em \RB{86.621}{8.685}{3.225}{%
\RG{0.000}{0.003}{10.9091}{rm}{100}\RG{5.935}{0.003}{10.9091}{rm}{105}\RG{8.902}{0.003}{10.9091}{rm}{109}\RG{19.669}{0.003}{10.9091}{rm}{107}\RG{25.025}{0.003}{10.9091}{rm}{101}\RG{29.771}{0.003}{10.9091}{rm}{114}\RG{35.735}{0.003}{10.9091}{mi}{39}\RG{45.763}{0.003}{10.9091}{sy}{21}\RG{57.272}{0.003}{10.9091}{sy}{100}\RG{62.126}{0.003}{10.9091}{mi}{110}\RG{68.508}{0.003}{10.9091}{mi}{61}\RG{73.854}{0.003}{10.9091}{mi}{109}\RG{83.186}{0.003}{10.9091}{sy}{101}%
}\hskip.1em\allowbreak{} 。

\section*{2015A09}
线性空间\allowbreak{}\hskip.1em \RB{8.715}{7.954}{0.726}{%
\RG{0.000}{0.003}{10.9091}{mi}{86}%
}\hskip.1em\allowbreak{} 中的一族向量\allowbreak{}\hskip.1em \RB{54.389}{8.685}{3.225}{%
\RG{0.000}{0.003}{10.9091}{mi}{66}\RG{11.935}{0.003}{10.9091}{rm}{61}\RG{23.514}{0.003}{10.9091}{sy}{102}\RG{28.968}{0.003}{10.9091}{mi}{101}\RG{33.919}{-1.634}{7.9701}{mi}{105}\RG{36.993}{0.003}{10.9091}{sy}{103}\RG{42.447}{-1.634}{7.9701}{mi}{105}\RG{45.069}{-1.634}{7.9701}{sy}{50}\RG{50.201}{-1.634}{7.9701}{mi}{73}%
}\hskip.1em\allowbreak{} 称为线性无关的，如果其中任意有限个向量都线性无关。若\allowbreak{}\hskip.1em \RB{8.605}{7.954}{0.500}{%
\RG{0.000}{0.003}{10.9091}{mi}{66}%
}\hskip.1em\allowbreak{} 线性无关，且\allowbreak{}\hskip.1em \RB{47.192}{8.685}{3.235}{%
\RG{0.000}{0.003}{10.9091}{mi}{86}\RG{11.911}{0.003}{10.9091}{rm}{61}\RG{23.528}{0.003}{10.9091}{mi}{76}\RG{30.796}{0.003}{10.9091}{rm}{40}\RG{34.950}{0.003}{10.9091}{mi}{66}\RG{43.615}{0.003}{10.9091}{rm}{41}%
}\hskip.1em\allowbreak{} ，即\allowbreak{}\hskip.1em \RB{8.715}{7.954}{0.726}{%
\RG{0.000}{0.003}{10.9091}{mi}{86}%
}\hskip.1em\allowbreak{} 中每个向量都能表示为\allowbreak{}\hskip.1em \RB{8.605}{7.954}{0.500}{%
\RG{0.000}{0.003}{10.9091}{mi}{66}%
}\hskip.1em\allowbreak{} 中有限个向量的线性组合，则称\allowbreak{}\hskip.1em \RB{8.605}{7.954}{0.500}{%
\RG{0.000}{0.003}{10.9091}{mi}{66}%
}\hskip.1em\allowbreak{} 为\allowbreak{}\hskip.1em \RB{8.715}{7.954}{0.726}{%
\RG{0.000}{0.003}{10.9091}{mi}{86}%
}\hskip.1em\allowbreak{} 的一组基。利用Zorn引理或选择公理可证明任一线性空间都存在一组基，本题承认这一结论。（1）证明：\allowbreak{}\hskip.1em \RB{21.163}{8.680}{3.229}{%
\RG{0.000}{-0.002}{10.9091}{mi}{75}\RG{9.871}{-0.002}{10.9091}{rm}{91}\RG{12.902}{-0.002}{10.9091}{mi}{120}\RG{18.972}{-0.002}{10.9091}{rm}{93}%
}\hskip.1em\allowbreak{} 的一组基为\allowbreak{}\hskip.1em \RB{97.481}{9.276}{3.229}{%
\RG{0.000}{-0.002}{10.9091}{mi}{66}\RG{11.695}{-0.002}{10.9091}{rm}{61}\RG{23.034}{-0.002}{10.9091}{sy}{102}\RG{28.488}{-0.002}{10.9091}{rm}{49}\RG{33.829}{-0.002}{10.9091}{mi}{59}\RG{38.607}{-0.002}{10.9091}{mi}{120}\RG{44.673}{-0.002}{10.9091}{mi}{59}\RG{49.462}{-0.002}{10.9091}{mi}{120}\RG{55.540}{3.957}{7.9701}{rm}{50}\RG{59.842}{-0.002}{10.9091}{mi}{59}\RG{64.620}{-0.002}{10.9091}{mi}{120}\RG{70.697}{3.957}{7.9701}{rm}{51}\RG{74.999}{-0.002}{10.9091}{mi}{59}\RG{79.777}{-0.002}{10.9091}{mi}{58}\RG{84.566}{-0.002}{10.9091}{mi}{58}\RG{89.355}{-0.002}{10.9091}{mi}{58}\RG{92.323}{-0.002}{10.9091}{sy}{103}%
}\hskip.1em\allowbreak{} 。（2）举例说明：有限维空间中关于可逆线性变换的不变子空间结论，对无限维空间一般不成立；即存在无限维线性空间\allowbreak{}\hskip.1em \RB{8.715}{7.954}{0.726}{%
\RG{0.000}{0.003}{10.9091}{mi}{86}%
}\hskip.1em\allowbreak{} 上的自同构\allowbreak{}\hskip.1em \RB{7.122}{5.314}{2.854}{%
\RG{0.000}{0.003}{10.9091}{mi}{39}%
}\hskip.1em\allowbreak{} 以及\allowbreak{}\hskip.1em \RB{7.122}{5.314}{2.854}{%
\RG{0.000}{0.003}{10.9091}{mi}{39}%
}\hskip.1em\allowbreak{} -不变子空间\allowbreak{}\hskip.1em \RB{11.693}{7.954}{0.726}{%
\RG{0.000}{0.003}{10.9091}{mi}{87}%
}\hskip.1em\allowbreak{} ，但\allowbreak{}\hskip.1em \RB{11.693}{7.954}{0.726}{%
\RG{0.000}{0.003}{10.9091}{mi}{87}%
}\hskip.1em\allowbreak{} 不是\allowbreak{}\hskip.1em \RB{16.724}{9.280}{2.854}{%
\RG{0.000}{0.003}{10.9091}{mi}{39}\RG{6.997}{3.962}{7.9701}{sy}{0}\RG{12.985}{3.962}{7.9701}{rm}{49}%
}\hskip.1em\allowbreak{} -不变子空间。

% Original page 2.
证明：（1）任取有限个不同的非负整数\allowbreak{}\hskip.1em \RB{39.197}{7.670}{2.603}{%
\RG{0.000}{0.003}{10.9091}{mi}{105}\RG{3.644}{-1.634}{7.9701}{rm}{49}\RG{7.945}{0.003}{10.9091}{mi}{59}\RG{12.724}{0.003}{10.9091}{mi}{58}\RG{17.513}{0.003}{10.9091}{mi}{58}\RG{22.302}{0.003}{10.9091}{mi}{58}\RG{27.080}{0.003}{10.9091}{mi}{59}\RG{31.869}{0.003}{10.9091}{mi}{105}\RG{35.516}{-1.634}{7.9701}{mi}{115}%
}\hskip.1em\allowbreak{} 。若

\[
\RB{104.577}{9.798}{2.608}{%
\RG{14.970}{4.502}{7.9701}{mi}{105}\RG{75.538}{4.502}{7.9701}{mi}{105}\RG{78.159}{3.596}{5.9776}{mi}{115}\RG{17.591}{3.495}{5.9776}{rm}{49}\RG{0.000}{-0.003}{10.9091}{mi}{99}\RG{8.900}{-0.003}{10.9091}{mi}{120}\RG{24.242}{-0.003}{10.9091}{rm}{43}\RG{34.974}{-0.003}{10.9091}{sy}{1}\RG{39.817}{-0.003}{10.9091}{sy}{1}\RG{44.672}{-0.003}{10.9091}{sy}{1}\RG{50.125}{-0.003}{10.9091}{rm}{43}\RG{60.858}{-0.003}{10.9091}{mi}{99}\RG{69.467}{-0.003}{10.9091}{mi}{120}\RG{85.271}{-0.003}{10.9091}{rm}{61}\RG{96.605}{-0.003}{10.9091}{rm}{48}\RG{101.950}{-0.003}{10.9091}{mi}{59}\RG{4.597}{-1.639}{7.9701}{rm}{49}\RG{65.455}{-1.639}{7.9701}{mi}{115}%
}
\]
则由多项式相等的定义，比较各次幂的系数，得到\allowbreak{}\hskip.1em \RB{78.726}{7.756}{2.207}{%
\RG{0.000}{0.002}{10.9091}{mi}{99}\RG{4.597}{-1.635}{7.9701}{rm}{49}\RG{11.930}{0.002}{10.9091}{rm}{61}\RG{23.268}{0.002}{10.9091}{sy}{1}\RG{28.112}{0.002}{10.9091}{sy}{1}\RG{32.966}{0.002}{10.9091}{sy}{1}\RG{39.025}{0.002}{10.9091}{rm}{61}\RG{50.365}{0.002}{10.9091}{mi}{99}\RG{54.962}{-1.635}{7.9701}{mi}{115}\RG{62.005}{0.002}{10.9091}{rm}{61}\RG{73.339}{0.002}{10.9091}{rm}{48}%
}\hskip.1em\allowbreak{} ，故\allowbreak{}\hskip.1em \RB{8.605}{7.953}{0.500}{%
\RG{0.000}{0.002}{10.9091}{mi}{66}%
}\hskip.1em\allowbreak{} 线性无关。每个多项式本来就是有限个单项式的线性组合，因此\allowbreak{}\hskip.1em \RB{59.970}{8.684}{3.236}{%
\RG{0.000}{0.002}{10.9091}{mi}{75}\RG{9.871}{0.002}{10.9091}{rm}{91}\RG{12.901}{0.002}{10.9091}{mi}{120}\RG{18.971}{0.002}{10.9091}{rm}{93}\RG{24.960}{0.002}{10.9091}{rm}{61}\RG{36.307}{0.002}{10.9091}{mi}{76}\RG{43.575}{0.002}{10.9091}{rm}{40}\RG{47.729}{0.002}{10.9091}{mi}{66}\RG{56.394}{0.002}{10.9091}{rm}{41}%
}\hskip.1em\allowbreak{} ，从而\allowbreak{}\hskip.1em \RB{8.605}{7.953}{0.500}{%
\RG{0.000}{0.002}{10.9091}{mi}{66}%
}\hskip.1em\allowbreak{} 是\allowbreak{}\hskip.1em \RB{21.163}{8.684}{3.225}{%
\RG{0.000}{0.002}{10.9091}{mi}{75}\RG{9.871}{0.002}{10.9091}{rm}{91}\RG{12.902}{0.002}{10.9091}{mi}{120}\RG{18.972}{0.002}{10.9091}{rm}{93}%
}\hskip.1em\allowbreak{} 的一组基。（2）取

\[
\RB{332.860}{16.814}{11.631}{%
\RG{80.422}{15.952}{9.0569}{ex}{40}\RG{246.760}{15.952}{9.0569}{ex}{41}\RG{92.876}{11.785}{7.9701}{mi}{115}\RG{88.115}{9.069}{9.0569}{ex}{88}\RG{52.796}{4.505}{7.9701}{sy}{0}\RG{58.785}{4.505}{7.9701}{rm}{49}\RG{119.745}{4.505}{7.9701}{mi}{107}\RG{0.000}{0.002}{10.9091}{mi}{86}\RG{11.631}{0.002}{10.9091}{rm}{61}\RG{22.970}{0.002}{10.9091}{mi}{75}\RG{32.841}{0.002}{10.9091}{rm}{91}\RG{35.871}{0.002}{10.9091}{mi}{120}\RG{41.936}{0.002}{10.9091}{mi}{59}\RG{46.715}{0.002}{10.9091}{mi}{120}\RG{63.086}{0.002}{10.9091}{rm}{93}\RG{69.075}{0.002}{10.9091}{rm}{61}\RG{103.413}{0.002}{10.9091}{mi}{97}\RG{113.675}{0.002}{10.9091}{mi}{120}\RG{127.429}{0.002}{10.9091}{rm}{58}\RG{136.983}{0.002}{10.9091}{mi}{114}\RG{141.499}{0.002}{10.9091}{mi}{59}\RG{146.277}{0.002}{10.9091}{mi}{115}\RG{154.347}{0.002}{10.9091}{sy}{50}\RG{164.650}{0.002}{10.9091}{bb}{90}\RG{171.923}{0.002}{10.9091}{mi}{59}\RG{180.257}{0.002}{10.9091}{mi}{114}\RG{188.405}{0.002}{10.9091}{sy}{20}\RG{199.921}{0.002}{10.9091}{mi}{115}\RG{204.950}{0.002}{10.9091}{mi}{59}\RG{213.285}{0.002}{10.9091}{mi}{97}\RG{226.585}{0.002}{10.9091}{sy}{50}\RG{236.888}{0.002}{10.9091}{mi}{75}\RG{255.874}{0.002}{10.9091}{mi}{59}\RG{282.470}{0.002}{10.9091}{mi}{87}\RG{297.075}{0.002}{10.9091}{rm}{61}\RG{308.414}{0.002}{10.9091}{mi}{75}\RG{318.285}{0.002}{10.9091}{rm}{91}\RG{321.315}{0.002}{10.9091}{mi}{120}\RG{327.385}{0.002}{10.9091}{rm}{93}\RG{330.353}{0.002}{10.9091}{mi}{58}\RG{109.020}{-1.635}{7.9701}{mi}{107}\RG{218.901}{-1.635}{7.9701}{mi}{107}\RG{87.718}{-11.058}{7.9701}{mi}{107}\RG{91.920}{-11.058}{7.9701}{rm}{61}\RG{97.907}{-11.058}{7.9701}{mi}{114}%
}
\]
这里\allowbreak{}\hskip.1em \RB{8.715}{7.949}{0.731}{%
\RG{0.000}{-0.002}{10.9091}{mi}{86}%
}\hskip.1em\allowbreak{} 是所有有限Laurent多项式构成的线性空间；它包含线性无关组\allowbreak{}\hskip.1em \RB{59.290}{9.275}{3.229}{%
\RG{0.000}{-0.002}{10.9091}{sy}{102}\RG{5.455}{-0.002}{10.9091}{rm}{49}\RG{10.795}{-0.002}{10.9091}{mi}{59}\RG{15.574}{-0.002}{10.9091}{mi}{120}\RG{21.639}{-0.002}{10.9091}{mi}{59}\RG{26.428}{-0.002}{10.9091}{mi}{120}\RG{32.506}{3.956}{7.9701}{rm}{50}\RG{36.808}{-0.002}{10.9091}{mi}{59}\RG{41.586}{-0.002}{10.9091}{mi}{58}\RG{46.375}{-0.002}{10.9091}{mi}{58}\RG{51.165}{-0.002}{10.9091}{mi}{58}\RG{54.132}{-0.002}{10.9091}{sy}{103}%
}\hskip.1em\allowbreak{} ，故为无限维空间。定义

\[
\RB{109.602}{8.680}{3.240}{%
\RG{0.000}{-0.002}{10.9091}{mi}{39}\RG{6.997}{-0.002}{10.9091}{rm}{40}\RG{11.152}{-0.002}{10.9091}{mi}{102}\RG{17.582}{-0.002}{10.9091}{rm}{41}\RG{24.760}{-0.002}{10.9091}{rm}{61}\RG{36.104}{-0.002}{10.9091}{mi}{120}\RG{42.169}{-0.002}{10.9091}{mi}{102}\RG{70.422}{-0.002}{10.9091}{rm}{40}\RG{74.576}{-0.002}{10.9091}{mi}{102}\RG{84.036}{-0.002}{10.9091}{sy}{50}\RG{94.339}{-0.002}{10.9091}{mi}{86}\RG{102.940}{-0.002}{10.9091}{rm}{41}\RG{107.095}{-0.002}{10.9091}{mi}{58}%
}
\]
乘以\allowbreak{}\hskip.1em \RB{6.107}{5.313}{0.607}{%
\RG{0.000}{0.002}{10.9091}{mi}{120}%
}\hskip.1em\allowbreak{} 与乘以\allowbreak{}\hskip.1em \RB{15.797}{9.279}{0.607}{%
\RG{0.000}{0.002}{10.9091}{mi}{120}\RG{6.070}{3.961}{7.9701}{sy}{0}\RG{12.058}{3.961}{7.9701}{rm}{49}%
}\hskip.1em\allowbreak{} 都将\allowbreak{}\hskip.1em \RB{8.715}{7.953}{0.727}{%
\RG{0.000}{0.002}{10.9091}{mi}{86}%
}\hskip.1em\allowbreak{} 映入自身，且互为逆映射，故\allowbreak{}\hskip.1em \RB{7.122}{5.313}{2.855}{%
\RG{0.000}{0.002}{10.9091}{mi}{39}%
}\hskip.1em\allowbreak{} 是\allowbreak{}\hskip.1em \RB{8.715}{7.953}{0.727}{%
\RG{0.000}{0.002}{10.9091}{mi}{86}%
}\hskip.1em\allowbreak{} 上的自同构，并且\allowbreak{}\hskip.1em \RB{43.262}{9.279}{3.236}{%
\RG{0.000}{0.002}{10.9091}{mi}{39}\RG{6.997}{3.961}{7.9701}{sy}{0}\RG{12.985}{3.961}{7.9701}{rm}{49}\RG{17.287}{0.002}{10.9091}{rm}{40}\RG{21.441}{0.002}{10.9091}{mi}{102}\RG{27.871}{0.002}{10.9091}{rm}{41}\RG{35.049}{0.002}{10.9091}{rm}{61}%
}\hskip.1em\allowbreak{}\allowbreak{}\hskip.1em \RB{22.816}{9.279}{2.724}{%
\RG{0.000}{0.002}{10.9091}{mi}{120}\RG{6.070}{3.961}{7.9701}{sy}{0}\RG{12.058}{3.961}{7.9701}{rm}{49}\RG{16.360}{0.002}{10.9091}{mi}{102}%
}\hskip.1em\allowbreak{} 。又有

\[
\RB{132.516}{8.684}{3.236}{%
\RG{0.000}{0.002}{10.9091}{mi}{39}\RG{6.997}{0.002}{10.9091}{rm}{40}\RG{11.151}{0.002}{10.9091}{mi}{87}\RG{22.719}{0.002}{10.9091}{rm}{41}\RG{29.897}{0.002}{10.9091}{rm}{61}\RG{41.242}{0.002}{10.9091}{mi}{120}\RG{47.307}{0.002}{10.9091}{mi}{75}\RG{57.183}{0.002}{10.9091}{rm}{91}\RG{60.214}{0.002}{10.9091}{mi}{120}\RG{66.284}{0.002}{10.9091}{rm}{93}\RG{72.281}{0.002}{10.9091}{sy}{18}\RG{83.796}{0.002}{10.9091}{mi}{75}\RG{93.667}{0.002}{10.9091}{rm}{91}\RG{96.697}{0.002}{10.9091}{mi}{120}\RG{102.767}{0.002}{10.9091}{rm}{93}\RG{108.756}{0.002}{10.9091}{rm}{61}\RG{120.104}{0.002}{10.9091}{mi}{87}\RG{129.889}{0.002}{10.9091}{mi}{59}%
}
\]
故\allowbreak{}\hskip.1em \RB{11.693}{7.947}{0.733}{%
\RG{0.000}{-0.004}{10.9091}{mi}{87}%
}\hskip.1em\allowbreak{} 是\allowbreak{}\hskip.1em \RB{7.122}{5.307}{2.860}{%
\RG{0.000}{-0.004}{10.9091}{mi}{39}%
}\hskip.1em\allowbreak{} -不变子空间；但\allowbreak{}\hskip.1em \RB{30.367}{7.947}{0.940}{%
\RG{0.000}{-0.004}{10.9091}{rm}{49}\RG{8.371}{-0.004}{10.9091}{sy}{50}\RG{18.675}{-0.004}{10.9091}{mi}{87}%
}\hskip.1em\allowbreak{} 而

\[
\RB{87.410}{9.824}{3.237}{%
\RG{6.997}{4.505}{7.9701}{sy}{0}\RG{12.985}{4.505}{7.9701}{rm}{49}\RG{51.375}{4.505}{7.9701}{sy}{0}\RG{57.363}{4.505}{7.9701}{rm}{49}\RG{0.000}{0.001}{10.9091}{mi}{39}\RG{17.287}{0.001}{10.9091}{rm}{40}\RG{21.444}{0.001}{10.9091}{rm}{49}\RG{26.789}{0.001}{10.9091}{rm}{41}\RG{33.967}{0.001}{10.9091}{rm}{61}\RG{45.305}{0.001}{10.9091}{mi}{120}\RG{64.695}{0.001}{10.9091}{sy}{50}\RG{65.964}{0.001}{10.9091}{mi}{61}\RG{74.997}{0.001}{10.9091}{mi}{87}\RG{84.783}{0.001}{10.9091}{mi}{59}%
}
\]
所以\allowbreak{}\hskip.1em \RB{11.693}{7.947}{0.733}{%
\RG{0.000}{-0.004}{10.9091}{mi}{87}%
}\hskip.1em\allowbreak{} 不是\allowbreak{}\hskip.1em \RB{16.724}{9.274}{2.860}{%
\RG{0.000}{-0.004}{10.9091}{mi}{39}\RG{6.997}{3.955}{7.9701}{sy}{0}\RG{12.985}{3.955}{7.9701}{rm}{49}%
}\hskip.1em\allowbreak{} -不变子空间。

小结：无限维空间中的线性组合仍只允许有限项。有限维情形可由限制映射的单射性推出满射性，而本例中\allowbreak{}\hskip.1em \RB{18.571}{8.678}{3.231}{%
\RG{0.000}{-0.004}{10.9091}{mi}{39}\RG{6.997}{-0.004}{10.9091}{sy}{106}\RG{10.028}{-1.640}{7.9701}{mi}{87}%
}\hskip.1em\allowbreak{} 单射但不满射，这正是原结论失效的原因。

\section*{2015A10}
设\allowbreak{}\hskip.1em \RB{8.715}{7.947}{0.733}{%
\RG{0.000}{-0.004}{10.9091}{mi}{86}%
}\hskip.1em\allowbreak{} 是数域\allowbreak{}\hskip.1em \RB{10.035}{7.947}{0.504}{%
\RG{0.000}{-0.004}{10.9091}{mi}{75}%
}\hskip.1em\allowbreak{} 上的\allowbreak{}\hskip.1em \RB{6.555}{5.307}{0.613}{%
\RG{0.000}{-0.004}{10.9091}{mi}{110}%
}\hskip.1em\allowbreak{} 维线性空间，\allowbreak{}\hskip.1em \RB{7.122}{5.307}{2.860}{%
\RG{0.000}{-0.004}{10.9091}{mi}{39}%
}\hskip.1em\allowbreak{} 是\allowbreak{}\hskip.1em \RB{8.715}{7.947}{0.733}{%
\RG{0.000}{-0.004}{10.9091}{mi}{86}%
}\hskip.1em\allowbreak{} 上的线性变换。证明下列条件等价：

\[
\RB{129.812}{53.241}{3.225}{%
\RT{0.000}{43.948}{10.9091}{（}\RT{10.909}{43.948}{10.9091}{1}\RT{16.255}{43.948}{10.9091}{）}\RG{30.982}{43.948}{10.9091}{mi}{86}\RG{42.614}{43.948}{10.9091}{rm}{61}\RG{53.948}{43.948}{10.9091}{rm}{107}\RG{59.305}{43.948}{10.9091}{rm}{101}\RG{64.050}{43.948}{10.9091}{rm}{114}\RG{70.012}{43.948}{10.9091}{mi}{39}\RG{79.434}{43.948}{10.9091}{rm}{43}\RG{90.168}{43.948}{10.9091}{rm}{73}\RG{94.019}{43.948}{10.9091}{rm}{109}\RG{104.792}{43.948}{10.9091}{mi}{39}\RT{111.789}{43.948}{10.9091}{。}\RT{0.000}{21.975}{10.9091}{（}\RT{10.909}{21.975}{10.9091}{2}\RT{16.255}{21.975}{10.9091}{）}\RG{30.982}{21.975}{10.9091}{mi}{86}\RG{42.614}{21.975}{10.9091}{rm}{61}\RG{53.948}{21.975}{10.9091}{rm}{107}\RG{59.305}{21.975}{10.9091}{rm}{101}\RG{64.050}{21.975}{10.9091}{rm}{114}\RG{70.012}{21.975}{10.9091}{mi}{39}\RG{79.434}{21.975}{10.9091}{sy}{8}\RG{90.343}{21.975}{10.9091}{rm}{73}\RG{94.194}{21.975}{10.9091}{rm}{109}\RG{104.968}{21.975}{10.9091}{mi}{39}\RT{111.965}{21.975}{10.9091}{。}\RT{0.000}{0.003}{10.9091}{（}\RT{10.909}{0.003}{10.9091}{3}\RT{16.255}{0.003}{10.9091}{）}\RG{30.982}{0.003}{10.9091}{rm}{107}\RG{36.327}{0.003}{10.9091}{rm}{101}\RG{41.073}{0.003}{10.9091}{rm}{114}\RG{47.043}{0.003}{10.9091}{mi}{39}\RG{56.464}{0.003}{10.9091}{sy}{92}\RG{66.161}{0.003}{10.9091}{rm}{73}\RG{70.012}{0.003}{10.9091}{rm}{109}\RG{80.786}{0.003}{10.9091}{mi}{39}\RG{90.815}{0.003}{10.9091}{rm}{61}\RG{102.153}{0.003}{10.9091}{sy}{102}\RG{107.607}{0.003}{10.9091}{rm}{48}\RG{112.948}{0.003}{10.9091}{sy}{103}\RT{118.403}{0.003}{10.9091}{。}%
}
\]
（4）\allowbreak{}\hskip.1em \RB{64.441}{9.279}{2.854}{%
\RG{0.000}{0.003}{10.9091}{rm}{107}\RG{5.345}{0.003}{10.9091}{rm}{101}\RG{10.091}{0.003}{10.9091}{rm}{114}\RG{16.061}{0.003}{10.9091}{mi}{39}\RG{26.088}{0.003}{10.9091}{rm}{61}\RG{37.423}{0.003}{10.9091}{rm}{107}\RG{42.779}{0.003}{10.9091}{rm}{101}\RG{47.525}{0.003}{10.9091}{rm}{114}\RG{53.487}{0.003}{10.9091}{mi}{39}\RG{60.485}{3.961}{7.9701}{rm}{50}%
}\hskip.1em\allowbreak{} ，或等价地，\allowbreak{}\hskip.1em \RB{103.788}{9.279}{2.854}{%
\RG{0.000}{0.003}{10.9091}{rm}{100}\RG{5.935}{0.003}{10.9091}{rm}{105}\RG{8.902}{0.003}{10.9091}{rm}{109}\RG{19.669}{0.003}{10.9091}{rm}{107}\RG{25.025}{0.003}{10.9091}{rm}{101}\RG{29.771}{0.003}{10.9091}{rm}{114}\RG{35.735}{0.003}{10.9091}{mi}{39}\RG{45.762}{0.003}{10.9091}{rm}{61}\RG{57.096}{0.003}{10.9091}{rm}{100}\RG{63.031}{0.003}{10.9091}{rm}{105}\RG{65.998}{0.003}{10.9091}{rm}{109}\RG{76.776}{0.003}{10.9091}{rm}{107}\RG{82.133}{0.003}{10.9091}{rm}{101}\RG{86.878}{0.003}{10.9091}{rm}{114}\RG{92.835}{0.003}{10.9091}{mi}{39}\RG{99.832}{3.961}{7.9701}{rm}{50}%
}\hskip.1em\allowbreak{} 。

（5）\allowbreak{}\hskip.1em \RB{61.572}{9.279}{2.854}{%
\RG{0.000}{0.003}{10.9091}{rm}{73}\RG{3.851}{0.003}{10.9091}{rm}{109}\RG{14.626}{0.003}{10.9091}{mi}{39}\RG{24.654}{0.003}{10.9091}{rm}{61}\RG{35.988}{0.003}{10.9091}{rm}{73}\RG{39.839}{0.003}{10.9091}{rm}{109}\RG{50.618}{0.003}{10.9091}{mi}{39}\RG{57.615}{3.961}{7.9701}{rm}{50}%
}\hskip.1em\allowbreak{} ，或等价地，\allowbreak{}\hskip.1em \RB{58.925}{9.279}{3.235}{%
\RG{0.000}{0.003}{10.9091}{mi}{114}\RG{5.110}{0.003}{10.9091}{rm}{40}\RG{9.265}{0.003}{10.9091}{mi}{39}\RG{16.262}{0.003}{10.9091}{rm}{41}\RG{23.440}{0.003}{10.9091}{rm}{61}\RG{34.785}{0.003}{10.9091}{mi}{114}\RG{39.895}{0.003}{10.9091}{rm}{40}\RG{44.049}{0.003}{10.9091}{mi}{39}\RG{51.046}{3.961}{7.9701}{rm}{50}\RG{55.348}{0.003}{10.9091}{rm}{41}%
}\hskip.1em\allowbreak{} 。

（6）\allowbreak{}\hskip.1em \RB{23.183}{8.074}{2.854}{%
\RG{0.000}{0.003}{10.9091}{rm}{107}\RG{5.345}{0.003}{10.9091}{rm}{101}\RG{10.091}{0.003}{10.9091}{rm}{114}\RG{16.061}{0.003}{10.9091}{mi}{39}%
}\hskip.1em\allowbreak{} 存在\allowbreak{}\hskip.1em \RB{7.122}{5.314}{2.854}{%
\RG{0.000}{0.003}{10.9091}{mi}{39}%
}\hskip.1em\allowbreak{} -不变的补空间，即存在\allowbreak{}\hskip.1em \RB{7.122}{5.314}{2.854}{%
\RG{0.000}{0.003}{10.9091}{mi}{39}%
}\hskip.1em\allowbreak{} -不变子空间\allowbreak{}\hskip.1em \RB{8.562}{7.954}{0.726}{%
\RG{0.000}{0.003}{10.9091}{mi}{85}%
}\hskip.1em\allowbreak{} ，使得\allowbreak{}\hskip.1em \RB{67.922}{8.074}{2.854}{%
\RG{0.000}{0.003}{10.9091}{mi}{86}\RG{11.631}{0.003}{10.9091}{rm}{61}\RG{22.965}{0.003}{10.9091}{rm}{107}\RG{28.322}{0.003}{10.9091}{rm}{101}\RG{33.067}{0.003}{10.9091}{rm}{114}\RG{39.030}{0.003}{10.9091}{mi}{39}\RG{48.451}{0.003}{10.9091}{sy}{8}\RG{59.360}{0.003}{10.9091}{mi}{85}%
}\hskip.1em\allowbreak{} 。

（7）\allowbreak{}\hskip.1em \RB{21.748}{7.953}{2.855}{%
\RG{0.000}{0.002}{10.9091}{rm}{73}\RG{3.851}{0.002}{10.9091}{rm}{109}\RG{14.626}{0.002}{10.9091}{mi}{39}%
}\hskip.1em\allowbreak{} 存在\allowbreak{}\hskip.1em \RB{7.122}{5.313}{2.855}{%
\RG{0.000}{0.002}{10.9091}{mi}{39}%
}\hskip.1em\allowbreak{} -不变的补空间，即存在\allowbreak{}\hskip.1em \RB{7.122}{5.313}{2.855}{%
\RG{0.000}{0.002}{10.9091}{mi}{39}%
}\hskip.1em\allowbreak{} -不变子空间\allowbreak{}\hskip.1em \RB{11.693}{7.953}{0.727}{%
\RG{0.000}{0.002}{10.9091}{mi}{87}%
}\hskip.1em\allowbreak{} ，使得\allowbreak{}\hskip.1em \RB{69.618}{7.953}{2.855}{%
\RG{0.000}{0.002}{10.9091}{mi}{86}\RG{11.631}{0.002}{10.9091}{rm}{61}\RG{22.965}{0.002}{10.9091}{rm}{73}\RG{26.816}{0.002}{10.9091}{rm}{109}\RG{37.595}{0.002}{10.9091}{mi}{39}\RG{47.016}{0.002}{10.9091}{sy}{8}\RG{57.925}{0.002}{10.9091}{mi}{87}%
}\hskip.1em\allowbreak{} 。

% Original page 3.
证明：记\allowbreak{}\hskip.1em \RB{47.244}{8.074}{2.854}{%
\RG{0.000}{0.003}{10.9091}{mi}{78}\RG{12.724}{0.003}{10.9091}{rm}{61}\RG{24.058}{0.003}{10.9091}{rm}{107}\RG{29.415}{0.003}{10.9091}{rm}{101}\RG{34.160}{0.003}{10.9091}{rm}{114}\RG{40.122}{0.003}{10.9091}{mi}{39}%
}\hskip.1em\allowbreak{} 、\allowbreak{}\hskip.1em \RB{44.337}{7.954}{2.854}{%
\RG{0.000}{0.003}{10.9091}{mi}{82}\RG{11.251}{0.003}{10.9091}{rm}{61}\RG{22.585}{0.003}{10.9091}{rm}{73}\RG{26.436}{0.003}{10.9091}{rm}{109}\RG{37.215}{0.003}{10.9091}{mi}{39}%
}\hskip.1em\allowbreak{} 。由秩与零度定理及维数公式，

\[
\RB{267.784}{8.679}{3.241}{%
\RG{0.000}{-0.003}{10.9091}{rm}{100}\RG{5.935}{-0.003}{10.9091}{rm}{105}\RG{8.902}{-0.003}{10.9091}{rm}{109}\RG{19.674}{-0.003}{10.9091}{mi}{78}\RG{31.791}{-0.003}{10.9091}{rm}{43}\RG{42.525}{-0.003}{10.9091}{rm}{100}\RG{48.460}{-0.003}{10.9091}{rm}{105}\RG{51.427}{-0.003}{10.9091}{rm}{109}\RG{62.197}{-0.003}{10.9091}{mi}{82}\RG{73.448}{-0.003}{10.9091}{rm}{61}\RG{84.786}{-0.003}{10.9091}{mi}{110}\RG{91.168}{-0.003}{10.9091}{mi}{59}\RG{117.766}{-0.003}{10.9091}{rm}{100}\RG{123.701}{-0.003}{10.9091}{rm}{105}\RG{126.668}{-0.003}{10.9091}{rm}{109}\RG{135.614}{-0.003}{10.9091}{rm}{40}\RG{139.775}{-0.003}{10.9091}{mi}{78}\RG{151.894}{-0.003}{10.9091}{rm}{43}\RG{162.625}{-0.003}{10.9091}{mi}{82}\RG{170.845}{-0.003}{10.9091}{rm}{41}\RG{178.024}{-0.003}{10.9091}{rm}{61}\RG{189.368}{-0.003}{10.9091}{mi}{110}\RG{198.169}{-0.003}{10.9091}{sy}{0}\RG{209.078}{-0.003}{10.9091}{rm}{100}\RG{215.013}{-0.003}{10.9091}{rm}{105}\RG{217.980}{-0.003}{10.9091}{rm}{109}\RG{226.925}{-0.003}{10.9091}{rm}{40}\RG{231.087}{-0.003}{10.9091}{mi}{78}\RG{243.205}{-0.003}{10.9091}{sy}{92}\RG{252.902}{-0.003}{10.9091}{mi}{82}\RG{261.123}{-0.003}{10.9091}{rm}{41}\RG{265.276}{-0.003}{10.9091}{mi}{58}%
}
\]
因此，\allowbreak{}\hskip.1em \RB{53.419}{7.953}{1.338}{%
\RG{0.000}{0.002}{10.9091}{mi}{78}\RG{11.751}{0.002}{10.9091}{rm}{43}\RG{22.115}{0.002}{10.9091}{mi}{82}\RG{33.366}{0.002}{10.9091}{rm}{61}\RG{44.705}{0.002}{10.9091}{mi}{86}%
}\hskip.1em\allowbreak{} 当且仅当\allowbreak{}\hskip.1em \RB{59.623}{8.684}{3.225}{%
\RG{0.000}{0.002}{10.9091}{mi}{78}\RG{11.750}{0.002}{10.9091}{sy}{92}\RG{21.080}{0.002}{10.9091}{mi}{82}\RG{32.331}{0.002}{10.9091}{rm}{61}\RG{43.669}{0.002}{10.9091}{sy}{102}\RG{49.124}{0.002}{10.9091}{rm}{48}\RG{54.465}{0.002}{10.9091}{sy}{103}%
}\hskip.1em\allowbreak{} ，而此时的和必为直和，故（1）、（2）、（3）等价。

若（3）成立，对任意\allowbreak{}\hskip.1em \RB{45.902}{9.275}{2.859}{%
\RG{0.000}{-0.003}{10.9091}{mi}{118}\RG{8.585}{-0.003}{10.9091}{sy}{50}\RG{18.889}{-0.003}{10.9091}{rm}{107}\RG{24.235}{-0.003}{10.9091}{rm}{101}\RG{28.980}{-0.003}{10.9091}{rm}{114}\RG{34.949}{-0.003}{10.9091}{mi}{39}\RG{41.946}{3.956}{7.9701}{rm}{50}%
}\hskip.1em\allowbreak{} ，有\allowbreak{}\hskip.1em \RB{91.554}{8.679}{3.241}{%
\RG{0.000}{-0.003}{10.9091}{mi}{39}\RG{6.997}{-0.003}{10.9091}{rm}{40}\RG{11.152}{-0.003}{10.9091}{mi}{118}\RG{16.707}{-0.003}{10.9091}{rm}{41}\RG{23.891}{-0.003}{10.9091}{sy}{50}\RG{34.195}{-0.003}{10.9091}{mi}{78}\RG{44.813}{-0.003}{10.9091}{sy}{92}\RG{53.011}{-0.003}{10.9091}{mi}{82}\RG{64.262}{-0.003}{10.9091}{rm}{61}\RG{75.600}{-0.003}{10.9091}{sy}{102}\RG{81.055}{-0.003}{10.9091}{rm}{48}\RG{86.395}{-0.003}{10.9091}{sy}{103}%
}\hskip.1em\allowbreak{} ，故\allowbreak{}\hskip.1em \RB{28.748}{7.948}{0.939}{%
\RG{0.000}{-0.003}{10.9091}{mi}{118}\RG{8.585}{-0.003}{10.9091}{sy}{50}\RG{18.888}{-0.003}{10.9091}{mi}{78}%
}\hskip.1em\allowbreak{} 。结合\allowbreak{}\hskip.1em \RB{64.617}{9.275}{2.859}{%
\RG{0.000}{-0.003}{10.9091}{rm}{107}\RG{5.345}{-0.003}{10.9091}{rm}{101}\RG{10.091}{-0.003}{10.9091}{rm}{114}\RG{16.061}{-0.003}{10.9091}{mi}{39}\RG{26.088}{-0.003}{10.9091}{sy}{18}\RG{37.604}{-0.003}{10.9091}{rm}{107}\RG{42.949}{-0.003}{10.9091}{rm}{101}\RG{47.695}{-0.003}{10.9091}{rm}{114}\RG{53.664}{-0.003}{10.9091}{mi}{39}\RG{60.661}{3.956}{7.9701}{rm}{50}%
}\hskip.1em\allowbreak{} ，可得（4）。反之，若（4）成立，任取\allowbreak{}\hskip.1em \RB{48.075}{7.948}{0.939}{%
\RG{0.000}{-0.003}{10.9091}{mi}{118}\RG{8.586}{-0.003}{10.9091}{sy}{50}\RG{18.889}{-0.003}{10.9091}{mi}{78}\RG{30.405}{-0.003}{10.9091}{sy}{92}\RG{39.502}{-0.003}{10.9091}{mi}{82}%
}\hskip.1em\allowbreak{} ，可写成\allowbreak{}\hskip.1em \RB{40.732}{8.679}{3.241}{%
\RG{0.000}{-0.003}{10.9091}{mi}{118}\RG{8.585}{-0.003}{10.9091}{rm}{61}\RG{19.925}{-0.003}{10.9091}{mi}{39}\RG{26.922}{-0.003}{10.9091}{rm}{40}\RG{31.075}{-0.003}{10.9091}{mi}{117}\RG{37.155}{-0.003}{10.9091}{rm}{41}%
}\hskip.1em\allowbreak{} ；于是\allowbreak{}\hskip.1em \RB{80.674}{9.275}{3.241}{%
\RG{0.000}{-0.003}{10.9091}{mi}{39}\RG{6.997}{3.956}{7.9701}{rm}{50}\RG{11.299}{-0.003}{10.9091}{rm}{40}\RG{15.454}{-0.003}{10.9091}{mi}{117}\RG{21.533}{-0.003}{10.9091}{rm}{41}\RG{28.711}{-0.003}{10.9091}{rm}{61}\RG{40.055}{-0.003}{10.9091}{mi}{39}\RG{47.053}{-0.003}{10.9091}{rm}{40}\RG{51.206}{-0.003}{10.9091}{mi}{118}\RG{56.763}{-0.003}{10.9091}{rm}{41}\RG{63.941}{-0.003}{10.9091}{rm}{61}\RG{75.286}{-0.003}{10.9091}{rm}{48}%
}\hskip.1em\allowbreak{} ，从而\allowbreak{}\hskip.1em \RB{84.324}{9.275}{2.859}{%
\RG{0.000}{-0.003}{10.9091}{mi}{117}\RG{9.110}{-0.003}{10.9091}{sy}{50}\RG{19.413}{-0.003}{10.9091}{rm}{107}\RG{24.758}{-0.003}{10.9091}{rm}{101}\RG{29.504}{-0.003}{10.9091}{rm}{114}\RG{35.474}{-0.003}{10.9091}{mi}{39}\RG{42.471}{3.956}{7.9701}{rm}{50}\RG{49.803}{-0.003}{10.9091}{rm}{61}\RG{61.137}{-0.003}{10.9091}{rm}{107}\RG{66.494}{-0.003}{10.9091}{rm}{101}\RG{71.239}{-0.003}{10.9091}{rm}{114}\RG{77.202}{-0.003}{10.9091}{mi}{39}%
}\hskip.1em\allowbreak{} ，故\allowbreak{}\hskip.1em \RB{25.307}{7.752}{0.732}{%
\RG{0.000}{-0.003}{10.9091}{mi}{118}\RG{8.585}{-0.003}{10.9091}{rm}{61}\RG{19.920}{-0.003}{10.9091}{rm}{48}%
}\hskip.1em\allowbreak{} 。因此（3）与（4）等价。

由于

\[
\RB{156.283}{9.825}{2.855}{%
\RG{60.661}{4.506}{7.9701}{rm}{50}\RG{113.189}{4.506}{7.9701}{rm}{50}\RG{0.000}{0.002}{10.9091}{rm}{107}\RG{5.345}{0.002}{10.9091}{rm}{101}\RG{10.091}{0.002}{10.9091}{rm}{114}\RG{16.060}{0.002}{10.9091}{mi}{39}\RG{26.088}{0.002}{10.9091}{sy}{18}\RG{37.603}{0.002}{10.9091}{rm}{107}\RG{42.948}{0.002}{10.9091}{rm}{101}\RG{47.694}{0.002}{10.9091}{rm}{114}\RG{53.664}{0.002}{10.9091}{mi}{39}\RG{64.963}{0.002}{10.9091}{mi}{59}\RG{91.566}{0.002}{10.9091}{rm}{73}\RG{95.417}{0.002}{10.9091}{rm}{109}\RG{106.192}{0.002}{10.9091}{mi}{39}\RG{120.522}{0.002}{10.9091}{sy}{18}\RG{132.037}{0.002}{10.9091}{rm}{73}\RG{135.888}{0.002}{10.9091}{rm}{109}\RG{146.663}{0.002}{10.9091}{mi}{39}\RG{153.655}{0.002}{10.9091}{mi}{59}%
}
\]
有限维性保证上述两个包含关系分别在两边维数相等时成为等式。再由

\[
\RB{232.285}{9.824}{3.237}{%
\RG{166.560}{4.505}{7.9701}{rm}{50}\RG{221.203}{4.505}{7.9701}{rm}{50}\RG{0.000}{0.001}{10.9091}{rm}{100}\RG{5.935}{0.001}{10.9091}{rm}{105}\RG{8.902}{0.001}{10.9091}{rm}{109}\RG{19.669}{0.001}{10.9091}{rm}{107}\RG{25.025}{0.001}{10.9091}{rm}{101}\RG{29.771}{0.001}{10.9091}{rm}{114}\RG{35.734}{0.001}{10.9091}{mi}{39}\RG{45.762}{0.001}{10.9091}{rm}{61}\RG{57.100}{0.001}{10.9091}{mi}{110}\RG{65.900}{0.001}{10.9091}{sy}{0}\RG{76.809}{0.001}{10.9091}{mi}{114}\RG{81.920}{0.001}{10.9091}{rm}{40}\RG{86.074}{0.001}{10.9091}{mi}{39}\RG{93.071}{0.001}{10.9091}{rm}{41}\RG{97.225}{0.001}{10.9091}{mi}{59}\RG{123.829}{0.001}{10.9091}{rm}{100}\RG{129.764}{0.001}{10.9091}{rm}{105}\RG{132.731}{0.001}{10.9091}{rm}{109}\RG{143.498}{0.001}{10.9091}{rm}{107}\RG{148.855}{0.001}{10.9091}{rm}{101}\RG{153.600}{0.001}{10.9091}{rm}{114}\RG{159.563}{0.001}{10.9091}{mi}{39}\RG{173.893}{0.001}{10.9091}{rm}{61}\RG{185.231}{0.001}{10.9091}{mi}{110}\RG{194.032}{0.001}{10.9091}{sy}{0}\RG{204.941}{0.001}{10.9091}{mi}{114}\RG{210.051}{0.001}{10.9091}{rm}{40}\RG{214.205}{0.001}{10.9091}{mi}{39}\RG{225.505}{0.001}{10.9091}{rm}{41}\RG{229.658}{0.001}{10.9091}{mi}{59}%
}
\]
可知（4）与（5）等价，同时证明了这两个条件中各自两种表述的等价性。

核\allowbreak{}\hskip.1em \RB{9.860}{7.952}{0.500}{%
\RG{0.000}{0.001}{10.9091}{mi}{78}%
}\hskip.1em\allowbreak{} 与像\allowbreak{}\hskip.1em \RB{8.573}{7.952}{0.728}{%
\RG{0.000}{0.001}{10.9091}{mi}{82}%
}\hskip.1em\allowbreak{} 都是\allowbreak{}\hskip.1em \RB{7.122}{5.312}{2.855}{%
\RG{0.000}{0.001}{10.9091}{mi}{39}%
}\hskip.1em\allowbreak{} -不变子空间。因此，若（2）成立，分别取\allowbreak{}\hskip.1em \RB{31.454}{7.952}{0.728}{%
\RG{0.000}{0.001}{10.9091}{mi}{85}\RG{11.483}{0.001}{10.9091}{rm}{61}\RG{22.881}{0.001}{10.9091}{mi}{82}%
}\hskip.1em\allowbreak{} 、\allowbreak{}\hskip.1em \RB{35.915}{7.952}{0.728}{%
\RG{0.000}{0.001}{10.9091}{mi}{87}\RG{14.657}{0.001}{10.9091}{rm}{61}\RG{26.055}{0.001}{10.9091}{mi}{78}%
}\hskip.1em\allowbreak{} ，即得（6）和（7）。

若（6）成立，则由\allowbreak{}\hskip.1em \RB{54.558}{7.947}{1.409}{%
\RG{0.000}{-0.004}{10.9091}{mi}{86}\RG{11.632}{-0.004}{10.9091}{rm}{61}\RG{22.970}{-0.004}{10.9091}{mi}{78}\RG{35.087}{-0.004}{10.9091}{sy}{8}\RG{45.996}{-0.004}{10.9091}{mi}{85}%
}\hskip.1em\allowbreak{} 及\allowbreak{}\hskip.1em \RB{8.562}{7.947}{0.733}{%
\RG{0.000}{-0.004}{10.9091}{mi}{85}%
}\hskip.1em\allowbreak{} 的不变性，

\[
\RB{108.947}{8.683}{3.237}{%
\RG{0.000}{0.001}{10.9091}{mi}{82}\RG{11.251}{0.001}{10.9091}{rm}{61}\RG{22.589}{0.001}{10.9091}{mi}{39}\RG{29.586}{0.001}{10.9091}{rm}{40}\RG{33.740}{0.001}{10.9091}{mi}{86}\RG{42.341}{0.001}{10.9091}{rm}{41}\RG{49.519}{0.001}{10.9091}{rm}{61}\RG{60.864}{0.001}{10.9091}{mi}{39}\RG{67.862}{0.001}{10.9091}{rm}{40}\RG{72.015}{0.001}{10.9091}{mi}{85}\RG{80.409}{0.001}{10.9091}{rm}{41}\RG{87.594}{0.001}{10.9091}{sy}{18}\RG{99.109}{0.001}{10.9091}{mi}{85}\RG{106.320}{0.001}{10.9091}{mi}{59}%
}
\]
从而\allowbreak{}\hskip.1em \RB{105.111}{8.678}{3.231}{%
\RG{0.000}{-0.004}{10.9091}{mi}{78}\RG{12.117}{-0.004}{10.9091}{sy}{92}\RG{21.815}{-0.004}{10.9091}{mi}{82}\RG{33.065}{-0.004}{10.9091}{sy}{18}\RG{44.580}{-0.004}{10.9091}{mi}{78}\RG{56.697}{-0.004}{10.9091}{sy}{92}\RG{66.395}{-0.004}{10.9091}{mi}{85}\RG{77.819}{-0.004}{10.9091}{rm}{61}\RG{89.157}{-0.004}{10.9091}{sy}{102}\RG{94.612}{-0.004}{10.9091}{rm}{48}\RG{99.953}{-0.004}{10.9091}{sy}{103}%
}\hskip.1em\allowbreak{} ，即（3）成立。

若（7）成立，则由\allowbreak{}\hskip.1em \RB{11.693}{7.952}{0.728}{%
\RG{0.000}{0.001}{10.9091}{mi}{87}%
}\hskip.1em\allowbreak{} 的不变性，

\[
\RB{106.575}{8.678}{3.242}{%
\RG{0.000}{-0.004}{10.9091}{mi}{39}\RG{6.997}{-0.004}{10.9091}{rm}{40}\RG{11.152}{-0.004}{10.9091}{mi}{87}\RG{22.719}{-0.004}{10.9091}{rm}{41}\RG{29.904}{-0.004}{10.9091}{sy}{18}\RG{41.419}{-0.004}{10.9091}{mi}{87}\RG{55.411}{-0.004}{10.9091}{sy}{92}\RG{65.108}{-0.004}{10.9091}{mi}{82}\RG{76.358}{-0.004}{10.9091}{rm}{61}\RG{87.697}{-0.004}{10.9091}{sy}{102}\RG{93.152}{-0.004}{10.9091}{rm}{48}\RG{98.493}{-0.004}{10.9091}{sy}{103}\RG{103.947}{-0.004}{10.9091}{mi}{59}%
}
\]
故\allowbreak{}\hskip.1em \RB{35.974}{7.951}{1.995}{%
\RG{0.000}{0.000}{10.9091}{mi}{87}\RG{14.598}{0.000}{10.9091}{sy}{18}\RG{26.114}{0.000}{10.9091}{mi}{78}%
}\hskip.1em\allowbreak{} ，于是\allowbreak{}\hskip.1em \RB{124.967}{7.951}{1.995}{%
\RG{0.000}{0.000}{10.9091}{mi}{86}\RG{11.632}{0.000}{10.9091}{rm}{61}\RG{22.970}{0.000}{10.9091}{mi}{82}\RG{33.615}{0.000}{10.9091}{sy}{8}\RG{44.524}{0.000}{10.9091}{mi}{87}\RG{59.122}{0.000}{10.9091}{sy}{18}\RG{70.637}{0.000}{10.9091}{mi}{82}\RG{81.282}{0.000}{10.9091}{rm}{43}\RG{92.014}{0.000}{10.9091}{mi}{78}\RG{104.737}{0.000}{10.9091}{sy}{18}\RG{116.253}{0.000}{10.9091}{mi}{86}%
}\hskip.1em\allowbreak{} ，得到（1）。综上，七个条件等价。

小结：这些条件共同刻画了核与像互为补空间的情形；此时\allowbreak{}\hskip.1em \RB{7.122}{5.311}{2.856}{%
\RG{0.000}{0.000}{10.9091}{mi}{39}%
}\hskip.1em\allowbreak{} 在核上为零，在像上的限制是自同构。有限维假设保证了“同维包含即相等”以及限制映射的单射与满射等价。

\section*{2015A11}
设\allowbreak{}\hskip.1em \RB{138.004}{8.682}{3.238}{%
\RG{0.000}{0.000}{10.9091}{mi}{102}\RG{5.265}{-1.636}{7.9701}{rm}{49}\RG{9.567}{0.000}{10.9091}{rm}{40}\RG{13.721}{0.000}{10.9091}{mi}{120}\RG{19.791}{0.000}{10.9091}{rm}{41}\RG{23.945}{0.000}{10.9091}{mi}{59}\RG{28.724}{0.000}{10.9091}{mi}{102}\RG{33.995}{-1.636}{7.9701}{rm}{50}\RG{38.297}{0.000}{10.9091}{rm}{40}\RG{42.452}{0.000}{10.9091}{mi}{120}\RG{48.522}{0.000}{10.9091}{rm}{41}\RG{52.675}{0.000}{10.9091}{mi}{59}\RG{57.454}{0.000}{10.9091}{mi}{58}\RG{62.243}{0.000}{10.9091}{mi}{58}\RG{67.032}{0.000}{10.9091}{mi}{58}\RG{71.810}{0.000}{10.9091}{mi}{59}\RG{76.599}{0.000}{10.9091}{mi}{102}\RG{81.867}{-1.636}{7.9701}{mi}{109}\RG{89.130}{0.000}{10.9091}{rm}{40}\RG{93.284}{0.000}{10.9091}{mi}{120}\RG{99.354}{0.000}{10.9091}{rm}{41}\RG{106.538}{0.000}{10.9091}{sy}{50}\RG{116.841}{0.000}{10.9091}{mi}{75}\RG{126.712}{0.000}{10.9091}{rm}{91}\RG{129.743}{0.000}{10.9091}{mi}{120}\RG{135.813}{0.000}{10.9091}{rm}{93}%
}\hskip.1em\allowbreak{} ，证明：

\[
\RB{282.685}{11.374}{33.996}{%
\RG{0.000}{10.512}{9.0569}{ex}{16}\RG{145.955}{10.512}{9.0569}{ex}{17}\RG{165.733}{10.512}{9.0569}{ex}{16}\RG{274.649}{10.512}{9.0569}{ex}{17}\RG{5.409}{-0.005}{10.9091}{rm}{40}\RG{9.563}{-0.005}{10.9091}{mi}{102}\RG{19.130}{-0.005}{10.9091}{rm}{40}\RG{23.285}{-0.005}{10.9091}{mi}{120}\RG{29.355}{-0.005}{10.9091}{rm}{41}\RG{33.508}{-0.005}{10.9091}{mi}{59}\RG{38.286}{-0.005}{10.9091}{mi}{102}\RG{47.861}{-0.005}{10.9091}{rm}{40}\RG{52.015}{-0.005}{10.9091}{mi}{120}\RG{58.085}{-0.005}{10.9091}{rm}{41}\RG{62.241}{-0.005}{10.9091}{rm}{41}\RG{66.393}{-0.005}{10.9091}{mi}{59}\RG{71.171}{-0.005}{10.9091}{mi}{102}\RG{80.745}{-0.005}{10.9091}{rm}{40}\RG{84.899}{-0.005}{10.9091}{mi}{120}\RG{90.969}{-0.005}{10.9091}{rm}{41}\RG{95.123}{-0.005}{10.9091}{mi}{59}\RG{99.901}{-0.005}{10.9091}{mi}{58}\RG{104.690}{-0.005}{10.9091}{mi}{58}\RG{109.479}{-0.005}{10.9091}{mi}{58}\RG{114.257}{-0.005}{10.9091}{mi}{59}\RG{119.046}{-0.005}{10.9091}{mi}{102}\RG{131.577}{-0.005}{10.9091}{rm}{40}\RG{135.731}{-0.005}{10.9091}{mi}{120}\RG{141.801}{-0.005}{10.9091}{rm}{41}\RG{154.394}{-0.005}{10.9091}{rm}{61}\RG{171.142}{-0.005}{10.9091}{mi}{102}\RG{180.709}{-0.005}{10.9091}{rm}{40}\RG{184.863}{-0.005}{10.9091}{mi}{120}\RG{190.933}{-0.005}{10.9091}{rm}{41}\RG{195.086}{-0.005}{10.9091}{mi}{59}\RG{199.865}{-0.005}{10.9091}{mi}{102}\RG{209.439}{-0.005}{10.9091}{rm}{40}\RG{213.593}{-0.005}{10.9091}{mi}{120}\RG{219.663}{-0.005}{10.9091}{rm}{41}\RG{223.817}{-0.005}{10.9091}{mi}{59}\RG{228.595}{-0.005}{10.9091}{mi}{58}\RG{233.385}{-0.005}{10.9091}{mi}{58}\RG{238.163}{-0.005}{10.9091}{mi}{58}\RG{242.952}{-0.005}{10.9091}{mi}{59}\RG{247.741}{-0.005}{10.9091}{mi}{102}\RG{260.271}{-0.005}{10.9091}{rm}{40}\RG{264.425}{-0.005}{10.9091}{mi}{120}\RG{270.495}{-0.005}{10.9091}{rm}{41}\RG{280.058}{-0.005}{10.9091}{mi}{59}\RG{14.828}{-1.641}{7.9701}{rm}{49}\RG{43.558}{-1.641}{7.9701}{rm}{50}\RG{76.443}{-1.641}{7.9701}{rm}{51}\RG{124.315}{-1.641}{7.9701}{mi}{109}\RG{176.406}{-1.641}{7.9701}{rm}{49}\RG{205.137}{-1.641}{7.9701}{rm}{50}\RG{253.008}{-1.641}{7.9701}{mi}{109}\RG{3.419}{-17.564}{9.0569}{ex}{104}\RG{145.932}{-17.564}{9.0569}{ex}{105}\RG{164.577}{-17.564}{9.0569}{ex}{104}\RG{272.362}{-17.564}{9.0569}{ex}{105}\RG{7.696}{-28.081}{10.9091}{rm}{91}\RG{10.726}{-28.081}{10.9091}{mi}{102}\RG{20.294}{-28.081}{10.9091}{rm}{40}\RG{24.448}{-28.081}{10.9091}{mi}{120}\RG{30.518}{-28.081}{10.9091}{rm}{41}\RG{34.672}{-28.081}{10.9091}{mi}{59}\RG{39.450}{-28.081}{10.9091}{mi}{102}\RG{49.025}{-28.081}{10.9091}{rm}{40}\RG{53.178}{-28.081}{10.9091}{mi}{120}\RG{59.248}{-28.081}{10.9091}{rm}{41}\RG{63.405}{-28.081}{10.9091}{rm}{93}\RG{66.369}{-28.081}{10.9091}{mi}{59}\RG{71.147}{-28.081}{10.9091}{mi}{102}\RG{80.722}{-28.081}{10.9091}{rm}{40}\RG{84.875}{-28.081}{10.9091}{mi}{120}\RG{90.945}{-28.081}{10.9091}{rm}{41}\RG{95.100}{-28.081}{10.9091}{mi}{59}\RG{99.878}{-28.081}{10.9091}{mi}{58}\RG{104.667}{-28.081}{10.9091}{mi}{58}\RG{109.446}{-28.081}{10.9091}{mi}{58}\RG{114.235}{-28.081}{10.9091}{mi}{59}\RG{119.024}{-28.081}{10.9091}{mi}{102}\RG{131.554}{-28.081}{10.9091}{rm}{40}\RG{135.708}{-28.081}{10.9091}{mi}{120}\RG{141.778}{-28.081}{10.9091}{rm}{41}\RG{153.239}{-28.081}{10.9091}{rm}{61}\RG{168.855}{-28.081}{10.9091}{mi}{102}\RG{178.422}{-28.081}{10.9091}{rm}{40}\RG{182.575}{-28.081}{10.9091}{mi}{120}\RG{188.645}{-28.081}{10.9091}{rm}{41}\RG{192.799}{-28.081}{10.9091}{mi}{59}\RG{197.577}{-28.081}{10.9091}{mi}{102}\RG{207.152}{-28.081}{10.9091}{rm}{40}\RG{211.305}{-28.081}{10.9091}{mi}{120}\RG{217.375}{-28.081}{10.9091}{rm}{41}\RG{221.530}{-28.081}{10.9091}{mi}{59}\RG{226.308}{-28.081}{10.9091}{mi}{58}\RG{231.097}{-28.081}{10.9091}{mi}{58}\RG{235.876}{-28.081}{10.9091}{mi}{58}\RG{240.665}{-28.081}{10.9091}{mi}{59}\RG{245.454}{-28.081}{10.9091}{mi}{102}\RG{257.984}{-28.081}{10.9091}{rm}{40}\RG{262.137}{-28.081}{10.9091}{mi}{120}\RG{268.207}{-28.081}{10.9091}{rm}{41}\RG{276.638}{-28.081}{10.9091}{mi}{59}\RG{15.992}{-29.717}{7.9701}{rm}{49}\RG{44.722}{-29.717}{7.9701}{rm}{50}\RG{76.420}{-29.717}{7.9701}{rm}{51}\RG{124.292}{-29.717}{7.9701}{mi}{109}\RG{174.119}{-29.717}{7.9701}{rm}{49}\RG{202.850}{-29.717}{7.9701}{rm}{50}\RG{250.721}{-29.717}{7.9701}{mi}{109}%
}
\]
其中小括号表示多项式的最大公因式，中括号表示多项式的最小公倍式。

% Original page 4.
证明：非零的最大公因式和最小公倍式均取首一多项式；约定全为零时最大公因式为\allowbreak{}\hskip.1em \RB{5.387}{7.757}{0.726}{%
\RG{0.000}{0.003}{10.9091}{rm}{48}%
}\hskip.1em\allowbreak{} ，含有零多项式时最小公倍式为\allowbreak{}\hskip.1em \RB{5.387}{7.757}{0.726}{%
\RG{0.000}{0.003}{10.9091}{rm}{48}%
}\hskip.1em\allowbreak{} 。整除关系按\allowbreak{}\hskip.1em \RB{86.295}{8.685}{3.225}{%
\RG{0.000}{0.003}{10.9091}{mi}{103}\RG{8.799}{0.003}{10.9091}{sy}{106}\RG{15.123}{0.003}{10.9091}{mi}{104}\RG{27.857}{0.003}{10.9091}{sy}{40}\RG{36.955}{0.003}{10.9091}{sy}{41}\RG{54.444}{0.003}{10.9091}{mi}{104}\RG{63.885}{0.003}{10.9091}{rm}{61}\RG{75.486}{0.003}{10.9091}{mi}{103}\RG{80.985}{0.003}{10.9091}{mi}{113}%
}\hskip.1em\allowbreak{} （某个\allowbreak{}\hskip.1em \RB{40.148}{8.685}{3.225}{%
\RG{0.000}{0.003}{10.9091}{mi}{113}\RG{8.420}{0.003}{10.9091}{sy}{50}\RG{18.986}{0.003}{10.9091}{mi}{75}\RG{28.857}{0.003}{10.9091}{rm}{91}\RG{31.887}{0.003}{10.9091}{mi}{120}\RG{37.957}{0.003}{10.9091}{rm}{93}%
}\hskip.1em\allowbreak{} ）理解，因而允许\allowbreak{}\hskip.1em \RB{19.819}{8.685}{3.225}{%
\RG{0.000}{0.003}{10.9091}{rm}{48}\RG{8.371}{0.003}{10.9091}{sy}{106}\RG{14.432}{0.003}{10.9091}{rm}{48}%
}\hskip.1em\allowbreak{} 。

以下略去自变量\allowbreak{}\hskip.1em \RB{6.107}{5.309}{0.611}{%
\RG{0.000}{-0.002}{10.9091}{mi}{120}%
}\hskip.1em\allowbreak{} ，令\allowbreak{}\hskip.1em \RB{51.569}{8.680}{3.240}{%
\RG{0.000}{-0.002}{10.9091}{mi}{100}\RG{8.581}{-0.002}{10.9091}{rm}{61}\RG{19.915}{-0.002}{10.9091}{rm}{40}\RG{24.074}{-0.002}{10.9091}{mi}{102}\RG{29.338}{-1.638}{7.9701}{rm}{49}\RG{33.641}{-0.002}{10.9091}{mi}{59}\RG{38.419}{-0.002}{10.9091}{mi}{102}\RG{43.691}{-1.638}{7.9701}{rm}{50}\RG{47.993}{-0.002}{10.9091}{rm}{41}%
}\hskip.1em\allowbreak{} 。由最大公因式的整除性质，对任意\allowbreak{}\hskip.1em \RB{40.645}{8.680}{3.229}{%
\RG{0.000}{-0.002}{10.9091}{mi}{104}\RG{9.179}{-0.002}{10.9091}{sy}{50}\RG{19.482}{-0.002}{10.9091}{mi}{75}\RG{29.353}{-0.002}{10.9091}{rm}{91}\RG{32.384}{-0.002}{10.9091}{mi}{120}\RG{38.454}{-0.002}{10.9091}{rm}{93}%
}\hskip.1em\allowbreak{} ，

\allowbreak{}\hskip.1em \RB{24.244}{8.686}{3.223}{%
\RG{0.000}{0.005}{10.9091}{mi}{104}\RG{9.179}{0.005}{10.9091}{sy}{106}\RG{15.240}{0.005}{10.9091}{mi}{102}\RG{20.505}{-1.632}{7.9701}{rm}{49}%
}\hskip.1em\allowbreak{} 且\allowbreak{}\hskip.1em \RB{95.985}{8.686}{3.223}{%
\RG{0.000}{0.005}{10.9091}{mi}{104}\RG{9.179}{0.005}{10.9091}{sy}{106}\RG{15.240}{0.005}{10.9091}{mi}{102}\RG{20.505}{-1.632}{7.9701}{rm}{50}\RG{38.746}{0.005}{10.9091}{sy}{40}\RG{47.845}{0.005}{10.9091}{sy}{41}\RG{72.685}{0.005}{10.9091}{mi}{104}\RG{81.865}{0.005}{10.9091}{sy}{106}\RG{87.925}{0.005}{10.9091}{mi}{100}\RG{93.478}{0.005}{10.9091}{mi}{58}%
}\hskip.1em\allowbreak{}

因此，\allowbreak{}\hskip.1em \RB{60.183}{8.183}{2.723}{%
\RG{0.000}{0.003}{10.9091}{mi}{102}\RG{5.265}{-1.634}{7.9701}{rm}{49}\RG{9.567}{0.003}{10.9091}{mi}{59}\RG{14.345}{0.003}{10.9091}{mi}{102}\RG{19.617}{-1.634}{7.9701}{rm}{50}\RG{23.920}{0.003}{10.9091}{mi}{59}\RG{28.698}{0.003}{10.9091}{mi}{58}\RG{33.487}{0.003}{10.9091}{mi}{58}\RG{38.265}{0.003}{10.9091}{mi}{58}\RG{43.055}{0.003}{10.9091}{mi}{59}\RG{47.844}{0.003}{10.9091}{mi}{102}\RG{53.111}{-1.634}{7.9701}{mi}{109}%
}\hskip.1em\allowbreak{} 与\allowbreak{}\hskip.1em \RB{56.166}{8.183}{2.723}{%
\RG{0.000}{0.003}{10.9091}{mi}{100}\RG{5.553}{0.003}{10.9091}{mi}{59}\RG{10.331}{0.003}{10.9091}{mi}{102}\RG{15.601}{-1.634}{7.9701}{rm}{51}\RG{19.903}{0.003}{10.9091}{mi}{59}\RG{24.681}{0.003}{10.9091}{mi}{58}\RG{29.470}{0.003}{10.9091}{mi}{58}\RG{34.259}{0.003}{10.9091}{mi}{58}\RG{39.037}{0.003}{10.9091}{mi}{59}\RG{43.826}{0.003}{10.9091}{mi}{102}\RG{49.095}{-1.634}{7.9701}{mi}{109}%
}\hskip.1em\allowbreak{} 的公因式完全相同。两组多项式遂有相同的首一最大公因式（全为零时均为\allowbreak{}\hskip.1em \RB{5.387}{7.757}{0.726}{%
\RG{0.000}{0.003}{10.9091}{rm}{48}%
}\hskip.1em\allowbreak{} ），第一个等式成立。

再令\allowbreak{}\hskip.1em \RB{47.992}{8.680}{3.229}{%
\RG{0.000}{-0.002}{10.9091}{mi}{96}\RG{7.513}{-0.002}{10.9091}{rm}{61}\RG{18.847}{-0.002}{10.9091}{rm}{91}\RG{21.882}{-0.002}{10.9091}{mi}{102}\RG{27.146}{-1.638}{7.9701}{rm}{49}\RG{31.449}{-0.002}{10.9091}{mi}{59}\RG{36.227}{-0.002}{10.9091}{mi}{102}\RG{41.499}{-1.638}{7.9701}{rm}{50}\RG{45.801}{-0.002}{10.9091}{rm}{93}%
}\hskip.1em\allowbreak{} 。由最小公倍式的整除性质，对任意\allowbreak{}\hskip.1em \RB{40.645}{8.680}{3.229}{%
\RG{0.000}{-0.002}{10.9091}{mi}{104}\RG{9.179}{-0.002}{10.9091}{sy}{50}\RG{19.483}{-0.002}{10.9091}{mi}{75}\RG{29.354}{-0.002}{10.9091}{rm}{91}\RG{32.384}{-0.002}{10.9091}{mi}{120}\RG{38.454}{-0.002}{10.9091}{rm}{93}%
}\hskip.1em\allowbreak{} ，

\allowbreak{}\hskip.1em \RB{24.995}{8.677}{3.232}{%
\RG{0.000}{-0.005}{10.9091}{mi}{102}\RG{5.265}{-1.641}{7.9701}{rm}{49}\RG{12.597}{-0.005}{10.9091}{sy}{106}\RG{18.658}{-0.005}{10.9091}{mi}{104}%
}\hskip.1em\allowbreak{} 且\allowbreak{}\hskip.1em \RB{94.919}{8.677}{3.232}{%
\RG{0.000}{-0.005}{10.9091}{mi}{102}\RG{5.265}{-1.641}{7.9701}{rm}{50}\RG{12.598}{-0.005}{10.9091}{sy}{106}\RG{18.658}{-0.005}{10.9091}{mi}{104}\RG{38.746}{-0.005}{10.9091}{sy}{40}\RG{47.845}{-0.005}{10.9091}{sy}{41}\RG{72.686}{-0.005}{10.9091}{mi}{96}\RG{80.199}{-0.005}{10.9091}{sy}{106}\RG{86.259}{-0.005}{10.9091}{mi}{104}\RG{92.412}{-0.005}{10.9091}{mi}{58}%
}\hskip.1em\allowbreak{}

因此，\allowbreak{}\hskip.1em \RB{60.183}{8.182}{2.724}{%
\RG{0.000}{0.002}{10.9091}{mi}{102}\RG{5.265}{-1.635}{7.9701}{rm}{49}\RG{9.567}{0.002}{10.9091}{mi}{59}\RG{14.345}{0.002}{10.9091}{mi}{102}\RG{19.617}{-1.635}{7.9701}{rm}{50}\RG{23.919}{0.002}{10.9091}{mi}{59}\RG{28.697}{0.002}{10.9091}{mi}{58}\RG{33.486}{0.002}{10.9091}{mi}{58}\RG{38.275}{0.002}{10.9091}{mi}{58}\RG{43.054}{0.002}{10.9091}{mi}{59}\RG{47.843}{0.002}{10.9091}{mi}{102}\RG{53.111}{-1.635}{7.9701}{mi}{109}%
}\hskip.1em\allowbreak{} 与\allowbreak{}\hskip.1em \RB{55.098}{8.182}{2.724}{%
\RG{0.000}{0.002}{10.9091}{mi}{96}\RG{4.484}{0.002}{10.9091}{mi}{59}\RG{9.262}{0.002}{10.9091}{mi}{102}\RG{14.533}{-1.635}{7.9701}{rm}{51}\RG{18.835}{0.002}{10.9091}{mi}{59}\RG{23.613}{0.002}{10.9091}{mi}{58}\RG{28.402}{0.002}{10.9091}{mi}{58}\RG{33.180}{0.002}{10.9091}{mi}{58}\RG{37.969}{0.002}{10.9091}{mi}{59}\RG{42.758}{0.002}{10.9091}{mi}{102}\RG{48.026}{-1.635}{7.9701}{mi}{109}%
}\hskip.1em\allowbreak{} 的公倍式完全相同。按上述规范化约定，它们的最小公倍式相同，第二个等式也成立。

小结：将前两个多项式合并为最大公因式或最小公倍式，分别不会改变整组多项式的公因式集合或公倍式集合。证明这类结合性，直接使用整除性质即可；首一规范化保证结论是严格相等。

\section*{2015A12}
设循环矩阵

\[
\RB{148.631}{41.394}{35.667}{%
\RG{22.376}{40.541}{9.0569}{ex}{48}\RG{142.851}{40.541}{9.0569}{ex}{49}\RG{35.631}{31.087}{10.9091}{mi}{97}\RG{60.337}{31.087}{10.9091}{mi}{97}\RG{79.715}{31.087}{10.9091}{mi}{97}\RG{98.681}{31.087}{10.9091}{sy}{1}\RG{103.525}{31.087}{10.9091}{sy}{1}\RG{108.379}{31.087}{10.9091}{sy}{1}\RG{127.202}{31.087}{10.9091}{mi}{97}\RG{41.238}{29.451}{7.9701}{rm}{49}\RG{65.945}{29.451}{7.9701}{rm}{50}\RG{85.322}{29.451}{7.9701}{rm}{51}\RG{132.809}{29.451}{7.9701}{mi}{110}\RG{142.851}{24.600}{9.0569}{ex}{67}\RG{22.376}{24.600}{9.0569}{ex}{66}\RG{22.376}{19.166}{9.0569}{ex}{66}\RG{142.851}{19.167}{9.0569}{ex}{67}\RG{35.220}{16.360}{10.9091}{mi}{97}\RG{60.337}{16.360}{10.9091}{mi}{97}\RG{79.715}{16.360}{10.9091}{mi}{97}\RG{98.681}{16.360}{10.9091}{sy}{1}\RG{103.525}{16.360}{10.9091}{sy}{1}\RG{108.379}{16.360}{10.9091}{sy}{1}\RG{122.283}{16.360}{10.9091}{mi}{97}\RG{40.827}{14.724}{7.9701}{mi}{110}\RG{65.945}{14.724}{7.9701}{rm}{49}\RG{85.322}{14.724}{7.9701}{rm}{50}\RG{127.890}{14.724}{7.9701}{mi}{110}\RG{132.561}{14.724}{7.9701}{sy}{0}\RG{138.549}{14.724}{7.9701}{rm}{49}\RG{22.376}{13.732}{9.0569}{ex}{66}\RG{142.851}{13.732}{9.0569}{ex}{67}\RG{22.376}{8.297}{9.0569}{ex}{66}\RG{142.851}{8.297}{9.0569}{ex}{67}\RG{22.376}{2.864}{9.0569}{ex}{66}\RG{142.851}{2.864}{9.0569}{ex}{67}\RG{30.301}{1.633}{10.9091}{mi}{97}\RG{59.926}{1.633}{10.9091}{mi}{97}\RG{79.715}{1.633}{10.9091}{mi}{97}\RG{98.681}{1.633}{10.9091}{sy}{1}\RG{103.525}{1.633}{10.9091}{sy}{1}\RG{108.379}{1.633}{10.9091}{sy}{1}\RG{122.283}{1.633}{10.9091}{mi}{97}\RG{0.000}{0.136}{10.9091}{mi}{65}\RG{11.038}{0.136}{10.9091}{rm}{61}\RG{35.908}{-0.004}{7.9701}{mi}{110}\RG{40.579}{-0.004}{7.9701}{sy}{0}\RG{46.567}{-0.004}{7.9701}{rm}{49}\RG{65.534}{-0.004}{7.9701}{mi}{110}\RG{85.322}{-0.004}{7.9701}{rm}{49}\RG{127.890}{-0.004}{7.9701}{mi}{110}\RG{132.561}{-0.004}{7.9701}{sy}{0}\RG{138.549}{-0.004}{7.9701}{rm}{50}\RG{22.376}{-2.571}{9.0569}{ex}{66}\RG{142.851}{-2.571}{9.0569}{ex}{67}\RG{22.376}{-8.005}{9.0569}{ex}{66}\RG{142.851}{-8.005}{9.0569}{ex}{67}\RT{39.102}{-9.278}{10.9091}{.}\RT{63.808}{-9.278}{10.9091}{.}\RT{83.186}{-9.278}{10.9091}{.}\RT{131.084}{-9.278}{10.9091}{.}\RT{39.102}{-12.901}{10.9091}{.}\RT{63.808}{-12.901}{10.9091}{.}\RT{83.186}{-12.901}{10.9091}{.}\RT{131.084}{-12.901}{10.9091}{.}\RG{22.376}{-13.439}{9.0569}{ex}{66}\RG{142.851}{-13.439}{9.0569}{ex}{67}\RT{39.102}{-16.524}{10.9091}{.}\RT{63.808}{-16.523}{10.9091}{.}\RT{83.186}{-16.523}{10.9091}{.}\RT{131.084}{-16.523}{10.9091}{.}\RG{22.376}{-19.235}{9.0569}{ex}{64}\RG{142.851}{-19.235}{9.0569}{ex}{65}\RG{35.631}{-31.251}{10.9091}{mi}{97}\RG{60.337}{-31.251}{10.9091}{mi}{97}\RG{79.715}{-31.251}{10.9091}{mi}{97}\RG{98.681}{-31.251}{10.9091}{sy}{1}\RG{103.525}{-31.251}{10.9091}{sy}{1}\RG{108.379}{-31.251}{10.9091}{sy}{1}\RG{127.613}{-31.251}{10.9091}{mi}{97}\RG{41.238}{-32.887}{7.9701}{rm}{50}\RG{65.945}{-32.887}{7.9701}{rm}{51}\RG{85.322}{-32.887}{7.9701}{rm}{52}\RG{133.220}{-32.887}{7.9701}{rm}{49}%
}
\]
是非异阵。求证：\allowbreak{}\hskip.1em \RB{17.734}{9.278}{0.500}{%
\RG{0.000}{0.001}{10.9091}{mi}{65}\RG{8.007}{3.959}{7.9701}{sy}{0}\RG{13.995}{3.959}{7.9701}{rm}{49}%
}\hskip.1em\allowbreak{} 也是循环矩阵。

证明：\allowbreak{}\hskip.1em \RB{25.714}{7.751}{0.613}{%
\RG{0.000}{-0.003}{10.9091}{mi}{110}\RG{9.406}{-0.003}{10.9091}{rm}{61}\RG{20.741}{-0.003}{10.9091}{rm}{49}%
}\hskip.1em\allowbreak{} 时结论显然。下设\allowbreak{}\hskip.1em \RB{26.222}{7.751}{1.998}{%
\RG{0.000}{-0.003}{10.9091}{mi}{110}\RG{9.406}{-0.003}{10.9091}{sy}{21}\RG{20.922}{-0.003}{10.9091}{rm}{50}%
}\hskip.1em\allowbreak{} ，令\allowbreak{}\hskip.1em \RB{43.064}{5.307}{2.609}{%
\RG{0.000}{-0.003}{10.9091}{mi}{101}\RG{4.951}{-1.640}{7.9701}{rm}{49}\RG{9.253}{-0.003}{10.9091}{mi}{59}\RG{14.031}{-0.003}{10.9091}{mi}{58}\RG{18.820}{-0.003}{10.9091}{mi}{58}\RG{23.598}{-0.003}{10.9091}{mi}{58}\RG{28.387}{-0.003}{10.9091}{mi}{59}\RG{33.176}{-0.003}{10.9091}{mi}{101}\RG{38.130}{-1.640}{7.9701}{mi}{110}%
}\hskip.1em\allowbreak{} 为标准单位列向量，取循环置换矩阵

\[
\RB{115.542}{8.685}{3.235}{%
\RG{0.000}{0.003}{10.9091}{mi}{80}\RG{11.399}{0.003}{10.9091}{rm}{61}\RG{22.734}{0.003}{10.9091}{rm}{40}\RG{26.892}{0.003}{10.9091}{mi}{101}\RG{36.966}{0.003}{10.9091}{mi}{59}\RG{41.745}{0.003}{10.9091}{mi}{101}\RG{51.005}{0.003}{10.9091}{mi}{59}\RG{55.783}{0.003}{10.9091}{mi}{101}\RG{65.042}{0.003}{10.9091}{mi}{59}\RG{69.820}{0.003}{10.9091}{mi}{58}\RG{74.609}{0.003}{10.9091}{mi}{58}\RG{79.398}{0.003}{10.9091}{mi}{58}\RG{84.176}{0.003}{10.9091}{mi}{59}\RG{88.965}{0.003}{10.9091}{mi}{101}\RG{108.880}{0.003}{10.9091}{rm}{41}\RG{113.035}{0.003}{10.9091}{mi}{58}\RG{31.842}{-1.633}{7.9701}{mi}{110}\RG{46.702}{-1.633}{7.9701}{rm}{49}\RG{60.740}{-1.633}{7.9701}{rm}{50}\RG{93.919}{-1.633}{7.9701}{mi}{110}\RG{98.591}{-1.633}{7.9701}{sy}{0}\RG{104.578}{-1.633}{7.9701}{rm}{49}%
}
\]
以下下标按模\allowbreak{}\hskip.1em \RB{6.555}{5.312}{0.608}{%
\RG{0.000}{0.001}{10.9091}{mi}{110}%
}\hskip.1em\allowbreak{} 循环理解，即下标\allowbreak{}\hskip.1em \RB{5.387}{7.756}{0.728}{%
\RG{0.000}{0.001}{10.9091}{rm}{48}%
}\hskip.1em\allowbreak{} 视为\allowbreak{}\hskip.1em \RB{6.555}{5.312}{0.608}{%
\RG{0.000}{0.001}{10.9091}{mi}{110}%
}\hskip.1em\allowbreak{} ，下标\allowbreak{}\hskip.1em \RB{23.226}{7.756}{1.339}{%
\RG{0.000}{0.001}{10.9091}{mi}{110}\RG{8.163}{0.001}{10.9091}{rm}{43}\RG{18.254}{0.001}{10.9091}{rm}{49}%
}\hskip.1em\allowbreak{} 视为\allowbreak{}\hskip.1em \RB{4.973}{7.756}{0.500}{%
\RG{0.000}{0.001}{10.9091}{rm}{49}%
}\hskip.1em\allowbreak{} 。任取\allowbreak{}\hskip.1em \RB{6.555}{5.312}{0.608}{%
\RG{0.000}{0.001}{10.9091}{mi}{110}%
}\hskip.1em\allowbreak{} 阶矩阵\allowbreak{}\hskip.1em \RB{44.496}{8.683}{3.613}{%
\RG{0.000}{0.001}{10.9091}{mi}{88}\RG{12.753}{0.001}{10.9091}{rm}{61}\RG{24.087}{0.001}{10.9091}{rm}{40}\RG{28.245}{0.001}{10.9091}{mi}{120}\RG{34.315}{-1.635}{7.9701}{mi}{105}\RG{36.938}{-1.635}{7.9701}{mi}{106}\RG{40.920}{0.001}{10.9091}{rm}{41}%
}\hskip.1em\allowbreak{} ，由矩阵乘法可得

\[
\RB{173.156}{8.683}{3.613}{%
\RG{0.000}{0.001}{10.9091}{rm}{40}\RG{4.154}{0.001}{10.9091}{mi}{88}\RG{13.874}{0.001}{10.9091}{mi}{80}\RG{22.245}{0.001}{10.9091}{rm}{41}\RG{36.035}{0.001}{10.9091}{rm}{61}\RG{47.373}{0.001}{10.9091}{mi}{120}\RG{72.023}{0.001}{10.9091}{mi}{59}\RG{98.626}{0.001}{10.9091}{rm}{40}\RG{102.780}{0.001}{10.9091}{mi}{80}\RG{111.147}{0.001}{10.9091}{mi}{88}\RG{120.872}{0.001}{10.9091}{rm}{41}\RG{134.661}{0.001}{10.9091}{rm}{61}\RG{145.999}{0.001}{10.9091}{mi}{120}\RG{170.649}{0.001}{10.9091}{mi}{58}\RG{26.399}{-1.635}{7.9701}{mi}{105}\RG{29.022}{-1.635}{7.9701}{mi}{106}\RG{53.443}{-1.635}{7.9701}{mi}{105}\RG{56.066}{-1.635}{7.9701}{mi}{59}\RG{58.203}{-1.635}{7.9701}{mi}{106}\RG{61.733}{-1.635}{7.9701}{sy}{0}\RG{67.721}{-1.635}{7.9701}{rm}{49}\RG{125.025}{-1.635}{7.9701}{mi}{105}\RG{127.648}{-1.635}{7.9701}{mi}{106}\RG{152.069}{-1.635}{7.9701}{mi}{105}\RG{154.690}{-1.635}{7.9701}{rm}{43}\RG{160.675}{-1.635}{7.9701}{rm}{49}\RG{164.527}{-1.635}{7.9701}{mi}{59}\RG{166.665}{-1.635}{7.9701}{mi}{106}%
}
\]
因而

\[
\RB{238.094}{8.677}{3.619}{%
\RG{0.000}{-0.004}{10.9091}{mi}{88}\RG{9.720}{-0.004}{10.9091}{mi}{80}\RG{21.122}{-0.004}{10.9091}{rm}{61}\RG{32.461}{-0.004}{10.9091}{mi}{80}\RG{40.817}{-0.004}{10.9091}{mi}{88}\RG{64.491}{-0.004}{10.9091}{sy}{40}\RG{73.589}{-0.004}{10.9091}{sy}{41}\RG{98.431}{-0.004}{10.9091}{mi}{120}\RG{126.111}{-0.004}{10.9091}{rm}{61}\RG{137.449}{-0.004}{10.9091}{mi}{120}\RG{173.009}{-0.004}{10.9091}{rm}{40}\RG{177.165}{-0.004}{10.9091}{rm}{49}\RG{185.534}{-0.004}{10.9091}{sy}{20}\RG{197.049}{-0.004}{10.9091}{mi}{105}\RG{200.693}{-0.004}{10.9091}{mi}{59}\RG{205.471}{-0.004}{10.9091}{mi}{106}\RG{213.542}{-0.004}{10.9091}{sy}{20}\RG{225.056}{-0.004}{10.9091}{mi}{110}\RG{231.433}{-0.004}{10.9091}{rm}{41}\RG{235.586}{-0.004}{10.9091}{mi}{58}\RG{104.501}{-1.641}{7.9701}{mi}{105}\RG{107.122}{-1.641}{7.9701}{rm}{43}\RG{113.107}{-1.641}{7.9701}{rm}{49}\RG{116.958}{-1.641}{7.9701}{mi}{59}\RG{119.096}{-1.641}{7.9701}{mi}{106}\RG{143.519}{-1.641}{7.9701}{mi}{105}\RG{146.142}{-1.641}{7.9701}{mi}{59}\RG{148.279}{-1.641}{7.9701}{mi}{106}\RG{151.810}{-1.641}{7.9701}{sy}{0}\RG{157.797}{-1.641}{7.9701}{rm}{49}%
}
\]
右端恰好表示后一行由前一行循环右移一位得到，即\allowbreak{}\hskip.1em \RB{9.642}{7.950}{0.501}{%
\RG{0.000}{-0.001}{10.9091}{mi}{88}%
}\hskip.1em\allowbreak{} 为题设形式的循环矩阵。于是\allowbreak{}\hskip.1em \RB{27.619}{8.288}{0.501}{%
\RG{0.000}{-0.001}{10.9091}{mi}{65}\RG{8.007}{-0.001}{10.9091}{mi}{80}\RG{19.406}{-0.001}{10.9091}{rm}{61}%
}\hskip.1em\allowbreak{}\allowbreak{}\hskip.1em \RB{16.569}{8.288}{0.501}{%
\RG{0.000}{-0.001}{10.9091}{mi}{80}\RG{8.367}{-0.001}{10.9091}{mi}{65}%
}\hskip.1em\allowbreak{} 。在此等式两边同时左乘、右乘\allowbreak{}\hskip.1em \RB{17.734}{9.276}{0.501}{%
\RG{0.000}{-0.001}{10.9091}{mi}{65}\RG{8.007}{3.957}{7.9701}{sy}{0}\RG{13.995}{3.957}{7.9701}{rm}{49}%
}\hskip.1em\allowbreak{} ，得到

\[
\RB{69.017}{9.819}{0.503}{%
\RG{16.376}{4.501}{7.9701}{sy}{0}\RG{22.365}{4.501}{7.9701}{rm}{49}\RG{49.043}{4.501}{7.9701}{sy}{0}\RG{55.030}{4.501}{7.9701}{rm}{49}\RG{0.000}{-0.003}{10.9091}{mi}{80}\RG{8.367}{-0.003}{10.9091}{mi}{65}\RG{29.696}{-0.003}{10.9091}{rm}{61}\RG{41.035}{-0.003}{10.9091}{mi}{65}\RG{59.332}{-0.003}{10.9091}{mi}{80}\RG{66.510}{-0.003}{10.9091}{mi}{58}%
}
\]
再由上述刻画，\allowbreak{}\hskip.1em \RB{17.734}{9.280}{0.500}{%
\RG{0.000}{0.004}{10.9091}{mi}{65}\RG{8.007}{3.962}{7.9701}{sy}{0}\RG{13.995}{3.962}{7.9701}{rm}{49}%
}\hskip.1em\allowbreak{} 也是循环矩阵。

小结：循环结构可用与固定循环置换矩阵交换来刻画，而可逆矩阵的逆保留这一交换关系。此法直接说明逆矩阵仍逐行循环右移，无须计算逆矩阵的各个元素。

% Original page 5.
\section*{2016S01}
设

\[
\RB{211.109}{9.820}{3.240}{%
\RG{41.246}{4.502}{7.9701}{mi}{110}\RG{86.165}{4.502}{7.9701}{mi}{110}\RG{90.836}{4.502}{7.9701}{sy}{0}\RG{96.824}{4.502}{7.9701}{rm}{49}\RG{0.000}{-0.002}{10.9091}{mi}{102}\RG{6.430}{-0.002}{10.9091}{rm}{40}\RG{10.584}{-0.002}{10.9091}{mi}{120}\RG{16.654}{-0.002}{10.9091}{rm}{41}\RG{23.832}{-0.002}{10.9091}{rm}{61}\RG{35.176}{-0.002}{10.9091}{mi}{120}\RG{48.795}{-0.002}{10.9091}{rm}{43}\RG{59.527}{-0.002}{10.9091}{mi}{97}\RG{80.095}{-0.002}{10.9091}{mi}{120}\RG{103.550}{-0.002}{10.9091}{rm}{43}\RG{114.283}{-0.002}{10.9091}{sy}{1}\RG{119.126}{-0.002}{10.9091}{sy}{1}\RG{123.981}{-0.002}{10.9091}{sy}{1}\RG{129.434}{-0.002}{10.9091}{rm}{43}\RG{140.166}{-0.002}{10.9091}{mi}{97}\RG{150.075}{-0.002}{10.9091}{mi}{120}\RG{158.570}{-0.002}{10.9091}{rm}{43}\RG{169.303}{-0.002}{10.9091}{mi}{97}\RG{182.242}{-0.002}{10.9091}{sy}{50}\RG{192.545}{-0.002}{10.9091}{bb}{90}\RG{199.818}{-0.002}{10.9091}{rm}{91}\RG{202.848}{-0.002}{10.9091}{mi}{120}\RG{208.918}{-0.002}{10.9091}{rm}{93}\RG{65.135}{-1.638}{7.9701}{mi}{110}\RG{69.805}{-1.638}{7.9701}{sy}{0}\RG{75.793}{-1.638}{7.9701}{rm}{49}\RG{145.774}{-1.638}{7.9701}{rm}{49}\RG{174.910}{-1.638}{7.9701}{rm}{48}%
}
\]
是首一多项式，满足\allowbreak{}\hskip.1em \RB{15.175}{8.684}{3.225}{%
\RG{0.000}{0.002}{10.9091}{sy}{106}\RG{3.031}{0.002}{10.9091}{mi}{97}\RG{8.638}{-1.635}{7.9701}{rm}{48}\RG{12.940}{0.002}{10.9091}{sy}{106}%
}\hskip.1em\allowbreak{} 是素数，且

\[
\RB{82.411}{17.107}{11.385}{%
\RG{48.835}{11.789}{7.9701}{mi}{110}\RG{53.507}{11.789}{7.9701}{sy}{0}\RG{59.495}{11.789}{7.9701}{rm}{49}\RG{49.549}{9.072}{9.0569}{ex}{88}\RG{0.000}{0.005}{10.9091}{sy}{106}\RG{3.030}{0.005}{10.9091}{mi}{97}\RG{12.939}{0.005}{10.9091}{sy}{106}\RG{19.000}{0.005}{10.9091}{mi}{62}\RG{30.338}{0.005}{10.9091}{rm}{49}\RG{38.105}{0.005}{10.9091}{rm}{43}\RG{65.162}{0.005}{10.9091}{sy}{106}\RG{68.192}{0.005}{10.9091}{mi}{97}\RG{76.874}{0.005}{10.9091}{sy}{106}\RG{79.904}{0.005}{10.9091}{mi}{58}\RG{8.637}{-1.632}{7.9701}{rm}{48}\RG{73.799}{-1.632}{7.9701}{mi}{105}\RG{49.861}{-10.813}{7.9701}{mi}{105}\RG{52.482}{-10.813}{7.9701}{rm}{61}\RG{58.467}{-10.813}{7.9701}{rm}{49}%
}
\]
证明：\allowbreak{}\hskip.1em \RB{20.230}{8.680}{3.240}{%
\RG{0.000}{-0.002}{10.9091}{mi}{102}\RG{6.430}{-0.002}{10.9091}{rm}{40}\RG{10.584}{-0.002}{10.9091}{mi}{120}\RG{16.654}{-0.002}{10.9091}{rm}{41}%
}\hskip.1em\allowbreak{} 在有理数域上不可约。注：原书将上述判别法称为Osada定理。

证明：先证明\allowbreak{}\hskip.1em \RB{6.456}{8.182}{2.724}{%
\RG{0.000}{0.002}{10.9091}{mi}{102}%
}\hskip.1em\allowbreak{} 的每个复根\allowbreak{}\hskip.1em \RB{5.289}{8.095}{2.724}{%
\RG{0.000}{0.002}{10.9091}{mi}{24}%
}\hskip.1em\allowbreak{} 都满足\allowbreak{}\hskip.1em \RB{30.595}{8.684}{3.225}{%
\RG{0.000}{0.002}{10.9091}{sy}{106}\RG{3.031}{0.002}{10.9091}{mi}{24}\RG{8.223}{0.002}{10.9091}{sy}{106}\RG{14.284}{0.002}{10.9091}{mi}{62}\RG{25.622}{0.002}{10.9091}{rm}{49}%
}\hskip.1em\allowbreak{} 。否则，由\allowbreak{}\hskip.1em \RB{39.687}{8.684}{3.236}{%
\RG{0.000}{0.002}{10.9091}{mi}{102}\RG{6.430}{0.002}{10.9091}{rm}{40}\RG{10.584}{0.002}{10.9091}{mi}{24}\RG{15.776}{0.002}{10.9091}{rm}{41}\RG{22.955}{0.002}{10.9091}{rm}{61}\RG{34.300}{0.002}{10.9091}{rm}{48}%
}\hskip.1em\allowbreak{} 及\allowbreak{}\hskip.1em \RB{30.771}{8.684}{3.225}{%
\RG{0.000}{0.002}{10.9091}{sy}{106}\RG{3.030}{0.002}{10.9091}{mi}{24}\RG{8.223}{0.002}{10.9091}{sy}{106}\RG{14.277}{0.002}{10.9091}{sy}{20}\RG{25.798}{0.002}{10.9091}{rm}{49}%
}\hskip.1em\allowbreak{} ，得到

\[
\RB{251.154}{17.107}{11.544}{%
\RG{30.339}{16.317}{9.0569}{ex}{12}\RG{90.105}{16.317}{9.0569}{ex}{12}\RG{56.831}{11.789}{7.9701}{mi}{110}\RG{61.502}{11.789}{7.9701}{sy}{0}\RG{67.490}{11.789}{7.9701}{rm}{49}\RG{137.203}{11.789}{7.9701}{mi}{110}\RG{141.874}{11.789}{7.9701}{sy}{0}\RG{147.862}{11.789}{7.9701}{rm}{49}\RG{217.459}{11.789}{7.9701}{mi}{110}\RG{222.130}{11.789}{7.9701}{sy}{0}\RG{228.117}{11.789}{7.9701}{rm}{49}\RG{30.339}{10.883}{9.0569}{ex}{12}\RG{90.105}{10.883}{9.0569}{ex}{12}\RG{57.544}{9.072}{9.0569}{ex}{88}\RG{137.915}{9.072}{9.0569}{ex}{88}\RG{218.172}{9.072}{9.0569}{ex}{88}\RG{30.339}{5.449}{9.0569}{ex}{12}\RG{90.105}{5.449}{9.0569}{ex}{12}\RG{38.550}{4.509}{7.9701}{mi}{110}\RG{87.031}{4.509}{7.9701}{mi}{105}\RG{118.922}{4.509}{7.9701}{mi}{110}\RG{181.342}{4.509}{7.9701}{mi}{105}\RG{0.000}{0.005}{10.9091}{sy}{106}\RG{3.031}{0.005}{10.9091}{mi}{97}\RG{12.940}{0.005}{10.9091}{sy}{106}\RG{19.001}{0.005}{10.9091}{rm}{61}\RG{30.339}{0.015}{9.0569}{ex}{12}\RG{33.358}{0.005}{10.9091}{mi}{24}\RG{46.098}{0.005}{10.9091}{rm}{43}\RG{73.157}{0.005}{10.9091}{mi}{97}\RG{81.838}{0.005}{10.9091}{mi}{24}\RG{90.105}{0.015}{9.0569}{ex}{12}\RG{96.154}{0.005}{10.9091}{sy}{20}\RG{107.663}{0.005}{10.9091}{sy}{106}\RG{110.699}{0.005}{10.9091}{mi}{24}\RG{115.892}{0.005}{10.9091}{sy}{106}\RG{126.470}{0.005}{10.9091}{rm}{43}\RG{153.529}{0.005}{10.9091}{sy}{106}\RG{156.559}{0.005}{10.9091}{mi}{97}\RG{165.240}{0.005}{10.9091}{sy}{106}\RG{170.084}{0.005}{10.9091}{sy}{106}\RG{173.119}{0.005}{10.9091}{mi}{24}\RG{178.312}{0.005}{10.9091}{sy}{106}\RG{187.446}{0.005}{10.9091}{sy}{20}\RG{198.961}{0.005}{10.9091}{rm}{49}\RG{206.728}{0.005}{10.9091}{rm}{43}\RG{233.785}{0.005}{10.9091}{sy}{106}\RG{236.815}{0.005}{10.9091}{mi}{97}\RG{245.496}{0.005}{10.9091}{sy}{106}\RG{248.526}{0.005}{10.9091}{mi}{59}\RG{8.638}{-1.632}{7.9701}{rm}{48}\RG{78.765}{-1.632}{7.9701}{mi}{105}\RG{162.166}{-1.632}{7.9701}{mi}{105}\RG{242.423}{-1.632}{7.9701}{mi}{105}\RG{30.339}{-5.420}{9.0569}{ex}{12}\RG{90.105}{-5.420}{9.0569}{ex}{12}\RG{57.856}{-10.814}{7.9701}{mi}{105}\RG{60.477}{-10.814}{7.9701}{rm}{61}\RG{66.462}{-10.814}{7.9701}{rm}{49}\RG{138.228}{-10.814}{7.9701}{mi}{105}\RG{140.849}{-10.814}{7.9701}{rm}{61}\RG{146.834}{-10.814}{7.9701}{rm}{49}\RG{218.484}{-10.814}{7.9701}{mi}{105}\RG{221.104}{-10.814}{7.9701}{rm}{61}\RG{227.089}{-10.814}{7.9701}{rm}{49}%
}
\]
与题设矛盾。

若\allowbreak{}\hskip.1em \RB{6.456}{8.179}{2.726}{%
\RG{0.000}{-0.001}{10.9091}{mi}{102}%
}\hskip.1em\allowbreak{} 在\allowbreak{}\hskip.1em \RB{8.627}{8.179}{2.454}{%
\RG{0.000}{-0.001}{10.9091}{bb}{81}%
}\hskip.1em\allowbreak{} 上可约，由Gauss引理及\allowbreak{}\hskip.1em \RB{6.456}{8.179}{2.726}{%
\RG{0.000}{-0.001}{10.9091}{mi}{102}%
}\hskip.1em\allowbreak{} 首一，可写为

\[
\RB{252.833}{8.685}{3.235}{%
\RG{0.000}{0.003}{10.9091}{mi}{102}\RG{6.430}{0.003}{10.9091}{rm}{40}\RG{10.584}{0.003}{10.9091}{mi}{120}\RG{16.654}{0.003}{10.9091}{rm}{41}\RG{23.832}{0.003}{10.9091}{rm}{61}\RG{35.176}{0.003}{10.9091}{mi}{103}\RG{40.683}{0.003}{10.9091}{rm}{40}\RG{44.836}{0.003}{10.9091}{mi}{120}\RG{50.907}{0.003}{10.9091}{rm}{41}\RG{55.061}{0.003}{10.9091}{mi}{104}\RG{61.210}{0.003}{10.9091}{rm}{40}\RG{65.364}{0.003}{10.9091}{mi}{120}\RG{71.434}{0.003}{10.9091}{rm}{41}\RG{75.587}{0.003}{10.9091}{mi}{59}\RG{102.184}{0.003}{10.9091}{mi}{103}\RG{107.693}{0.003}{10.9091}{mi}{59}\RG{112.482}{0.003}{10.9091}{mi}{104}\RG{121.662}{0.003}{10.9091}{sy}{50}\RG{131.965}{0.003}{10.9091}{bb}{90}\RG{139.237}{0.003}{10.9091}{rm}{91}\RG{142.268}{0.003}{10.9091}{mi}{120}\RG{148.338}{0.003}{10.9091}{rm}{93}\RG{151.305}{0.003}{10.9091}{mi}{59}\RG{177.909}{0.003}{10.9091}{rm}{100}\RG{183.844}{0.003}{10.9091}{rm}{101}\RG{188.589}{0.003}{10.9091}{rm}{103}\RG{195.894}{0.003}{10.9091}{mi}{103}\RG{201.392}{0.003}{10.9091}{mi}{59}\RG{206.185}{0.003}{10.9091}{rm}{100}\RG{212.120}{0.003}{10.9091}{rm}{101}\RG{216.865}{0.003}{10.9091}{rm}{103}\RG{224.170}{0.003}{10.9091}{mi}{104}\RG{233.349}{0.003}{10.9091}{sy}{21}\RG{244.865}{0.003}{10.9091}{rm}{49}\RG{250.205}{0.003}{10.9091}{mi}{59}%
}
\]
其中\allowbreak{}\hskip.1em \RB{16.624}{8.069}{2.727}{%
\RG{0.000}{-0.002}{10.9091}{mi}{103}\RG{5.498}{-0.002}{10.9091}{mi}{59}\RG{10.287}{-0.002}{10.9091}{mi}{104}%
}\hskip.1em\allowbreak{} 均首一。由

\[
\RB{82.435}{8.680}{3.240}{%
\RG{0.000}{-0.002}{10.9091}{sy}{106}\RG{3.030}{-0.002}{10.9091}{mi}{103}\RG{8.536}{-0.002}{10.9091}{rm}{40}\RG{12.693}{-0.002}{10.9091}{rm}{48}\RG{18.038}{-0.002}{10.9091}{rm}{41}\RG{22.185}{-0.002}{10.9091}{sy}{106}\RG{27.029}{-0.002}{10.9091}{sy}{106}\RG{30.064}{-0.002}{10.9091}{mi}{104}\RG{36.213}{-0.002}{10.9091}{rm}{40}\RG{40.369}{-0.002}{10.9091}{rm}{48}\RG{45.715}{-0.002}{10.9091}{rm}{41}\RG{49.862}{-0.002}{10.9091}{sy}{106}\RG{55.923}{-0.002}{10.9091}{rm}{61}\RG{67.261}{-0.002}{10.9091}{sy}{106}\RG{70.291}{-0.002}{10.9091}{mi}{97}\RG{80.200}{-0.002}{10.9091}{sy}{106}\RG{75.898}{-1.638}{7.9701}{rm}{48}%
}
\]
为素数，必有一个因子的常数项绝对值为\allowbreak{}\hskip.1em \RB{4.973}{7.757}{0.500}{%
\RG{0.000}{0.003}{10.9091}{rm}{49}%
}\hskip.1em\allowbreak{} ，不妨设\allowbreak{}\hskip.1em \RB{44.553}{8.685}{3.235}{%
\RG{0.000}{0.003}{10.9091}{sy}{106}\RG{3.030}{0.003}{10.9091}{mi}{103}\RG{8.536}{0.003}{10.9091}{rm}{40}\RG{12.693}{0.003}{10.9091}{rm}{48}\RG{18.038}{0.003}{10.9091}{rm}{41}\RG{22.185}{0.003}{10.9091}{sy}{106}\RG{28.245}{0.003}{10.9091}{rm}{61}\RG{39.580}{0.003}{10.9091}{rm}{49}%
}\hskip.1em\allowbreak{} 。设\allowbreak{}\hskip.1em \RB{5.595}{5.314}{2.723}{%
\RG{0.000}{0.003}{10.9091}{mi}{103}%
}\hskip.1em\allowbreak{} 的全部复根（按重数计）为\allowbreak{}\hskip.1em \RB{41.240}{8.095}{2.723}{%
\RG{0.000}{0.003}{10.9091}{mi}{24}\RG{4.665}{-1.634}{7.9701}{rm}{49}\RG{8.967}{0.003}{10.9091}{mi}{59}\RG{13.745}{0.003}{10.9091}{mi}{58}\RG{18.535}{0.003}{10.9091}{mi}{58}\RG{23.324}{0.003}{10.9091}{mi}{58}\RG{28.102}{0.003}{10.9091}{mi}{59}\RG{32.891}{0.003}{10.9091}{mi}{24}\RG{37.559}{-1.634}{7.9701}{mi}{115}%
}\hskip.1em\allowbreak{} ，其中\allowbreak{}\hskip.1em \RB{62.410}{8.074}{2.734}{%
\RG{0.000}{0.003}{10.9091}{mi}{115}\RG{8.062}{0.003}{10.9091}{rm}{61}\RG{19.396}{0.003}{10.9091}{rm}{100}\RG{25.331}{0.003}{10.9091}{rm}{101}\RG{30.076}{0.003}{10.9091}{rm}{103}\RG{37.385}{0.003}{10.9091}{mi}{103}\RG{45.922}{0.003}{10.9091}{sy}{21}\RG{57.437}{0.003}{10.9091}{rm}{49}%
}\hskip.1em\allowbreak{} 。它们都是\allowbreak{}\hskip.1em \RB{6.456}{8.183}{2.723}{%
\RG{0.000}{0.003}{10.9091}{mi}{102}%
}\hskip.1em\allowbreak{} 的根，故由Vieta公式，

\[
\RB{146.255}{15.480}{12.795}{%
\RG{98.926}{11.785}{7.9701}{mi}{115}\RG{94.920}{9.068}{9.0569}{ex}{89}\RG{0.000}{0.001}{10.9091}{sy}{106}\RG{3.030}{0.001}{10.9091}{mi}{103}\RG{8.536}{0.001}{10.9091}{rm}{40}\RG{12.693}{0.001}{10.9091}{rm}{48}\RG{18.038}{0.001}{10.9091}{rm}{41}\RG{22.185}{0.001}{10.9091}{sy}{106}\RG{28.245}{0.001}{10.9091}{rm}{61}\RG{39.584}{0.001}{10.9091}{sy}{106}\RG{42.615}{0.001}{10.9091}{mi}{24}\RG{53.400}{0.001}{10.9091}{sy}{1}\RG{58.244}{0.001}{10.9091}{sy}{1}\RG{63.098}{0.001}{10.9091}{sy}{1}\RG{67.945}{0.001}{10.9091}{mi}{24}\RG{76.624}{0.001}{10.9091}{sy}{106}\RG{82.684}{0.001}{10.9091}{rm}{61}\RG{109.208}{0.001}{10.9091}{sy}{106}\RG{112.238}{0.001}{10.9091}{mi}{24}\RG{120.887}{0.001}{10.9091}{sy}{106}\RG{126.948}{0.001}{10.9091}{mi}{62}\RG{138.286}{0.001}{10.9091}{rm}{49}\RG{143.627}{0.001}{10.9091}{mi}{59}\RG{47.280}{-1.635}{7.9701}{rm}{49}\RG{72.611}{-1.635}{7.9701}{mi}{115}\RG{116.904}{-1.635}{7.9701}{mi}{106}\RG{94.023}{-10.817}{7.9701}{mi}{106}\RG{97.554}{-10.817}{7.9701}{rm}{61}\RG{103.538}{-10.817}{7.9701}{rm}{49}%
}
\]
矛盾。因此\allowbreak{}\hskip.1em \RB{20.230}{8.678}{3.242}{%
\RG{0.000}{-0.004}{10.9091}{mi}{102}\RG{6.430}{-0.004}{10.9091}{rm}{40}\RG{10.584}{-0.004}{10.9091}{mi}{120}\RG{16.654}{-0.004}{10.9091}{rm}{41}%
}\hskip.1em\allowbreak{} 在\allowbreak{}\hskip.1em \RB{8.627}{8.176}{2.456}{%
\RG{0.000}{-0.004}{10.9091}{bb}{81}%
}\hskip.1em\allowbreak{} 上不可约。

小结：系数不等式保证全部根均在单位圆外；常数项的素性则迫使非平凡整系数分解中出现常数项绝对值为\allowbreak{}\hskip.1em \RB{4.973}{7.751}{0.504}{%
\RG{0.000}{-0.004}{10.9091}{rm}{49}%
}\hskip.1em\allowbreak{} 的因子，两者由Vieta公式发生矛盾。

% Original page 6.
\section*{2016S02}
（1）设\allowbreak{}\hskip.1em \RB{7.122}{5.309}{2.858}{%
\RG{0.000}{-0.002}{10.9091}{mi}{39}%
}\hskip.1em\allowbreak{} 是\allowbreak{}\hskip.1em \RB{6.555}{5.309}{0.611}{%
\RG{0.000}{-0.002}{10.9091}{mi}{110}%
}\hskip.1em\allowbreak{} 维线性空间\allowbreak{}\hskip.1em \RB{8.715}{7.949}{0.731}{%
\RG{0.000}{-0.002}{10.9091}{mi}{86}%
}\hskip.1em\allowbreak{} 上的线性变换，且

\[
\RB{112.873}{7.953}{2.604}{%
\RG{0.000}{0.002}{10.9091}{mi}{86}\RG{11.631}{0.002}{10.9091}{rm}{61}\RG{22.970}{0.002}{10.9091}{mi}{86}\RG{35.924}{0.002}{10.9091}{sy}{8}\RG{46.833}{0.002}{10.9091}{mi}{86}\RG{59.785}{0.002}{10.9091}{sy}{8}\RG{70.695}{0.002}{10.9091}{sy}{1}\RG{75.538}{0.002}{10.9091}{sy}{1}\RG{80.393}{0.002}{10.9091}{sy}{1}\RG{85.847}{0.002}{10.9091}{sy}{8}\RG{96.755}{0.002}{10.9091}{mi}{86}\RG{110.245}{0.002}{10.9091}{mi}{59}\RG{29.197}{-1.635}{7.9701}{rm}{49}\RG{53.060}{-1.635}{7.9701}{rm}{50}\RG{102.983}{-1.635}{7.9701}{mi}{109}%
}
\]
其中每个\allowbreak{}\hskip.1em \RB{9.024}{7.947}{2.212}{%
\RG{0.000}{-0.004}{10.9091}{mi}{86}\RG{6.227}{-1.640}{7.9701}{mi}{105}%
}\hskip.1em\allowbreak{} 都是\allowbreak{}\hskip.1em \RB{7.122}{5.307}{2.860}{%
\RG{0.000}{-0.004}{10.9091}{mi}{39}%
}\hskip.1em\allowbreak{} -不变子空间。设\allowbreak{}\hskip.1em \RB{10.274}{8.067}{2.292}{%
\RG{0.000}{-0.004}{10.9091}{mi}{21}\RG{6.231}{-1.640}{7.9701}{rm}{48}%
}\hskip.1em\allowbreak{} 是\allowbreak{}\hskip.1em \RB{7.122}{5.307}{2.860}{%
\RG{0.000}{-0.004}{10.9091}{mi}{39}%
}\hskip.1em\allowbreak{} 的特征值，记

\[
\RB{214.362}{8.682}{3.535}{%
\RG{0.000}{0.000}{10.9091}{mi}{86}\RG{13.559}{0.000}{10.9091}{rm}{61}\RG{24.897}{0.000}{10.9091}{sy}{102}\RG{30.352}{0.000}{10.9091}{mi}{118}\RG{38.938}{0.000}{10.9091}{sy}{50}\RG{49.241}{0.000}{10.9091}{mi}{86}\RG{60.872}{0.000}{10.9091}{rm}{58}\RG{66.870}{0.000}{10.9091}{mi}{39}\RG{73.867}{0.000}{10.9091}{rm}{40}\RG{78.021}{0.000}{10.9091}{mi}{118}\RG{83.576}{0.000}{10.9091}{rm}{41}\RG{90.755}{0.000}{10.9091}{rm}{61}\RG{102.099}{0.000}{10.9091}{mi}{21}\RG{112.632}{0.000}{10.9091}{mi}{118}\RG{118.188}{0.000}{10.9091}{sy}{103}\RG{123.643}{0.000}{10.9091}{mi}{59}\RG{150.239}{0.000}{10.9091}{mi}{86}\RG{168.565}{0.000}{10.9091}{rm}{61}\RG{179.904}{0.000}{10.9091}{mi}{86}\RG{191.628}{0.000}{10.9091}{sy}{92}\RG{201.325}{0.000}{10.9091}{mi}{86}\RG{211.855}{0.000}{10.9091}{mi}{58}\RG{6.227}{-1.636}{7.9701}{rm}{48}\RG{108.330}{-1.636}{7.9701}{rm}{48}\RG{156.473}{-1.636}{7.9701}{mi}{105}\RG{159.096}{-1.636}{7.9701}{mi}{59}\RG{161.233}{-1.636}{7.9701}{rm}{48}\RG{186.131}{-1.636}{7.9701}{mi}{105}\RG{207.553}{-1.636}{7.9701}{rm}{48}%
}
\]
证明：

\[
\RB{132.644}{7.955}{3.531}{%
\RG{0.000}{0.004}{10.9091}{mi}{86}\RG{13.559}{0.004}{10.9091}{rm}{61}\RG{24.898}{0.004}{10.9091}{mi}{86}\RG{43.839}{0.004}{10.9091}{sy}{8}\RG{54.748}{0.004}{10.9091}{mi}{86}\RG{73.689}{0.004}{10.9091}{sy}{8}\RG{84.598}{0.004}{10.9091}{sy}{1}\RG{89.442}{0.004}{10.9091}{sy}{1}\RG{94.296}{0.004}{10.9091}{sy}{1}\RG{99.751}{0.004}{10.9091}{sy}{8}\RG{110.659}{0.004}{10.9091}{mi}{86}\RG{130.136}{0.004}{10.9091}{mi}{58}\RG{6.227}{-1.633}{7.9701}{rm}{48}\RG{31.125}{-1.633}{7.9701}{rm}{49}\RG{34.975}{-1.633}{7.9701}{mi}{59}\RG{37.113}{-1.633}{7.9701}{rm}{48}\RG{60.975}{-1.633}{7.9701}{rm}{50}\RG{64.825}{-1.633}{7.9701}{mi}{59}\RG{66.963}{-1.633}{7.9701}{rm}{48}\RG{116.886}{-1.633}{7.9701}{mi}{109}\RG{123.697}{-1.633}{7.9701}{mi}{59}\RG{125.835}{-1.633}{7.9701}{rm}{48}%
}
\]
（2）设\allowbreak{}\hskip.1em \RB{6.555}{5.307}{0.613}{%
\RG{0.000}{-0.004}{10.9091}{mi}{110}%
}\hskip.1em\allowbreak{} 阶方阵

\[
\RB{121.640}{8.678}{3.242}{%
\RG{0.000}{-0.004}{10.9091}{mi}{65}\RG{11.038}{-0.004}{10.9091}{rm}{61}\RG{22.373}{-0.004}{10.9091}{rm}{100}\RG{28.307}{-0.004}{10.9091}{rm}{105}\RG{31.275}{-0.004}{10.9091}{rm}{97}\RG{36.620}{-0.004}{10.9091}{rm}{103}\RG{42.107}{-0.004}{10.9091}{rm}{40}\RG{46.258}{-0.004}{10.9091}{mi}{65}\RG{58.567}{-0.004}{10.9091}{mi}{59}\RG{63.345}{-0.004}{10.9091}{mi}{65}\RG{75.662}{-0.004}{10.9091}{mi}{59}\RG{80.440}{-0.004}{10.9091}{mi}{58}\RG{85.229}{-0.004}{10.9091}{mi}{58}\RG{90.018}{-0.004}{10.9091}{mi}{58}\RG{94.796}{-0.004}{10.9091}{mi}{59}\RG{99.585}{-0.004}{10.9091}{mi}{65}\RG{114.858}{-0.004}{10.9091}{rm}{41}\RG{119.013}{-0.004}{10.9091}{mi}{59}\RG{54.265}{-1.640}{7.9701}{rm}{49}\RG{71.360}{-1.640}{7.9701}{rm}{50}\RG{107.596}{-1.640}{7.9701}{mi}{109}%
}
\]
其中\allowbreak{}\hskip.1em \RB{10.804}{8.288}{2.210}{%
\RG{0.000}{-0.001}{10.9091}{mi}{65}\RG{8.007}{-1.637}{7.9701}{mi}{105}%
}\hskip.1em\allowbreak{} 是\allowbreak{}\hskip.1em \RB{9.173}{5.310}{2.210}{%
\RG{0.000}{-0.001}{10.9091}{mi}{110}\RG{6.376}{-1.637}{7.9701}{mi}{105}%
}\hskip.1em\allowbreak{} 阶方阵。任取\allowbreak{}\hskip.1em \RB{10.805}{8.288}{2.210}{%
\RG{0.000}{-0.001}{10.9091}{mi}{65}\RG{8.008}{-1.637}{7.9701}{mi}{105}%
}\hskip.1em\allowbreak{} 的特征值\allowbreak{}\hskip.1em \RB{9.028}{8.070}{2.210}{%
\RG{0.000}{-0.001}{10.9091}{mi}{21}\RG{6.231}{-1.637}{7.9701}{mi}{105}%
}\hskip.1em\allowbreak{} 及对应的特征向量\allowbreak{}\hskip.1em \RB{58.048}{8.310}{2.857}{%
\RG{0.000}{-0.001}{10.9091}{rm}{48}\RG{8.372}{-0.001}{10.9091}{sy}{54}\RG{8.372}{-0.001}{10.9091}{rm}{61}\RG{19.710}{-0.001}{10.9091}{mi}{11}\RG{26.504}{-1.637}{7.9701}{mi}{105}\RG{32.608}{-0.001}{10.9091}{sy}{50}\RG{42.911}{-0.001}{10.9091}{bb}{67}\RG{50.790}{3.957}{7.9701}{mi}{110}\RG{55.461}{2.853}{5.9776}{mi}{105}%
}\hskip.1em\allowbreak{} 。证明：在\allowbreak{}\hskip.1em \RB{9.590}{5.310}{2.210}{%
\RG{0.000}{-0.001}{10.9091}{mi}{11}\RG{6.794}{-1.637}{7.9701}{mi}{105}%
}\hskip.1em\allowbreak{}的上下添加适当多的零，所得非零向量\allowbreak{}\hskip.1em \RB{36.014}{8.179}{2.210}{%
\RG{0.000}{-0.001}{10.9091}{mi}{11}\RG{1.213}{0.796}{9.0569}{ex}{101}\RG{6.795}{-1.637}{7.9701}{mi}{105}\RG{12.898}{-0.001}{10.9091}{sy}{50}\RG{23.202}{-0.001}{10.9091}{bb}{67}\RG{31.080}{3.957}{7.9701}{mi}{110}%
}\hskip.1em\allowbreak{} 满足

\[
\RB{53.924}{8.292}{2.206}{%
\RG{9.220}{0.800}{9.0569}{ex}{101}\RG{42.761}{0.800}{9.0569}{ex}{101}\RG{0.000}{0.003}{10.9091}{mi}{65}\RG{8.007}{0.003}{10.9091}{mi}{11}\RG{20.905}{0.003}{10.9091}{rm}{61}\RG{32.244}{0.003}{10.9091}{mi}{21}\RG{41.549}{0.003}{10.9091}{mi}{11}\RG{51.416}{0.003}{10.9091}{mi}{58}\RG{14.801}{-1.634}{7.9701}{mi}{105}\RG{38.475}{-1.634}{7.9701}{mi}{105}\RG{48.343}{-1.634}{7.9701}{mi}{105}%
}
\]
称\allowbreak{}\hskip.1em \RB{9.591}{7.834}{2.212}{%
\RG{0.000}{-0.003}{10.9091}{mi}{11}\RG{1.213}{0.795}{9.0569}{ex}{101}\RG{6.795}{-1.639}{7.9701}{mi}{105}%
}\hskip.1em\allowbreak{} 为\allowbreak{}\hskip.1em \RB{9.591}{5.308}{2.212}{%
\RG{0.000}{-0.003}{10.9091}{mi}{11}\RG{6.795}{-1.639}{7.9701}{mi}{105}%
}\hskip.1em\allowbreak{} 的延拓。（3）在（2）的假设下，任取\allowbreak{}\hskip.1em \RB{8.202}{8.290}{0.500}{%
\RG{0.000}{0.001}{10.9091}{mi}{65}%
}\hskip.1em\allowbreak{} 的特征值\allowbreak{}\hskip.1em \RB{10.274}{8.072}{2.287}{%
\RG{0.000}{0.001}{10.9091}{mi}{21}\RG{6.231}{-1.635}{7.9701}{rm}{48}%
}\hskip.1em\allowbreak{} 。设\allowbreak{}\hskip.1em \RB{10.274}{8.072}{2.287}{%
\RG{0.000}{0.001}{10.9091}{mi}{21}\RG{6.231}{-1.635}{7.9701}{rm}{48}%
}\hskip.1em\allowbreak{} 是\allowbreak{}\hskip.1em \RB{52.954}{8.290}{3.141}{%
\RG{0.000}{0.001}{10.9091}{mi}{65}\RG{8.008}{-1.635}{7.9701}{mi}{105}\RG{10.629}{-2.641}{5.9776}{rm}{49}\RG{14.855}{0.001}{10.9091}{mi}{59}\RG{19.634}{0.001}{10.9091}{mi}{58}\RG{24.423}{0.001}{10.9091}{mi}{58}\RG{29.201}{0.001}{10.9091}{mi}{58}\RG{33.990}{0.001}{10.9091}{mi}{59}\RG{38.779}{0.001}{10.9091}{mi}{65}\RG{46.789}{-1.635}{7.9701}{mi}{105}\RG{49.411}{-2.541}{5.9776}{mi}{114}%
}\hskip.1em\allowbreak{} 的特征值，但不是其他块\allowbreak{}\hskip.1em \RB{11.500}{8.290}{3.613}{%
\RG{0.000}{0.001}{10.9091}{mi}{65}\RG{8.007}{-1.635}{7.9701}{mi}{106}%
}\hskip.1em\allowbreak{} 的特征值。证明：分别取\allowbreak{}\hskip.1em \RB{14.400}{8.290}{3.467}{%
\RG{0.000}{0.001}{10.9091}{mi}{65}\RG{8.007}{-1.635}{7.9701}{mi}{105}\RG{10.628}{-2.913}{5.9776}{mi}{107}%
}\hskip.1em\allowbreak{} 关于\allowbreak{}\hskip.1em \RB{10.274}{8.072}{2.288}{%
\RG{0.000}{0.001}{10.9091}{mi}{21}\RG{6.231}{-1.635}{7.9701}{rm}{48}%
}\hskip.1em\allowbreak{} 的特征子空间的一组基，将这些基向量延拓后取并集，即得\allowbreak{}\hskip.1em \RB{8.202}{8.290}{0.500}{%
\RG{0.000}{0.001}{10.9091}{mi}{65}%
}\hskip.1em\allowbreak{} 关于\allowbreak{}\hskip.1em \RB{10.275}{8.072}{2.287}{%
\RG{0.000}{0.001}{10.9091}{mi}{21}\RG{6.232}{-1.635}{7.9701}{rm}{48}%
}\hskip.1em\allowbreak{} 的特征子空间的一组基。

证明：（1）因为\allowbreak{}\hskip.1em \RB{38.858}{7.947}{3.538}{%
\RG{0.000}{-0.004}{10.9091}{mi}{86}\RG{6.227}{-1.640}{7.9701}{mi}{105}\RG{8.850}{-1.640}{7.9701}{mi}{59}\RG{10.986}{-1.640}{7.9701}{rm}{48}\RG{18.319}{-0.004}{10.9091}{sy}{18}\RG{29.835}{-0.004}{10.9091}{mi}{86}\RG{36.061}{-1.640}{7.9701}{mi}{105}%
}\hskip.1em\allowbreak{} ，而\allowbreak{}\hskip.1em \RB{61.954}{7.947}{2.213}{%
\RG{0.000}{-0.004}{10.9091}{mi}{86}\RG{6.227}{-1.640}{7.9701}{rm}{49}\RG{12.725}{-0.004}{10.9091}{rm}{43}\RG{23.228}{-0.004}{10.9091}{sy}{1}\RG{28.072}{-0.004}{10.9091}{sy}{1}\RG{32.926}{-0.004}{10.9091}{sy}{1}\RG{38.151}{-0.004}{10.9091}{rm}{43}\RG{48.655}{-0.004}{10.9091}{mi}{86}\RG{54.883}{-1.640}{7.9701}{mi}{109}%
}\hskip.1em\allowbreak{} 是直和，所以\allowbreak{}\hskip.1em \RB{73.861}{7.947}{3.538}{%
\RG{0.000}{-0.004}{10.9091}{mi}{86}\RG{6.227}{-1.640}{7.9701}{rm}{49}\RG{10.076}{-1.640}{7.9701}{mi}{59}\RG{12.215}{-1.640}{7.9701}{rm}{48}\RG{18.713}{-0.004}{10.9091}{rm}{43}\RG{29.216}{-0.004}{10.9091}{sy}{1}\RG{34.060}{-0.004}{10.9091}{sy}{1}\RG{38.915}{-0.004}{10.9091}{sy}{1}\RG{44.139}{-0.004}{10.9091}{rm}{43}\RG{54.643}{-0.004}{10.9091}{mi}{86}\RG{60.870}{-1.640}{7.9701}{mi}{109}\RG{67.681}{-1.640}{7.9701}{mi}{59}\RG{69.818}{-1.640}{7.9701}{rm}{48}%
}\hskip.1em\allowbreak{} 也是直和。又因\allowbreak{}\hskip.1em \RB{40.105}{7.947}{3.538}{%
\RG{0.000}{-0.004}{10.9091}{mi}{86}\RG{6.227}{-1.640}{7.9701}{mi}{105}\RG{8.850}{-1.640}{7.9701}{mi}{59}\RG{10.987}{-1.640}{7.9701}{rm}{48}\RG{18.319}{-0.004}{10.9091}{sy}{18}\RG{29.835}{-0.004}{10.9091}{mi}{86}\RG{36.062}{-1.640}{7.9701}{rm}{48}%
}\hskip.1em\allowbreak{} ，有

\[
\RB{102.970}{7.947}{3.538}{%
\RG{0.000}{-0.004}{10.9091}{mi}{86}\RG{18.942}{-0.004}{10.9091}{sy}{8}\RG{29.851}{-0.004}{10.9091}{sy}{1}\RG{34.695}{-0.004}{10.9091}{sy}{1}\RG{39.549}{-0.004}{10.9091}{sy}{1}\RG{45.004}{-0.004}{10.9091}{sy}{8}\RG{55.911}{-0.004}{10.9091}{mi}{86}\RG{78.419}{-0.004}{10.9091}{sy}{18}\RG{89.934}{-0.004}{10.9091}{mi}{86}\RG{100.463}{-0.004}{10.9091}{mi}{58}\RG{6.227}{-1.640}{7.9701}{rm}{49}\RG{10.076}{-1.640}{7.9701}{mi}{59}\RG{12.215}{-1.640}{7.9701}{rm}{48}\RG{62.138}{-1.640}{7.9701}{mi}{109}\RG{68.949}{-1.640}{7.9701}{mi}{59}\RG{71.086}{-1.640}{7.9701}{rm}{48}\RG{96.161}{-1.640}{7.9701}{rm}{48}%
}
\]
反之，任取\allowbreak{}\hskip.1em \RB{29.159}{7.949}{2.290}{%
\RG{0.000}{-0.002}{10.9091}{mi}{118}\RG{8.586}{-0.002}{10.9091}{sy}{50}\RG{18.889}{-0.002}{10.9091}{mi}{86}\RG{25.116}{-1.638}{7.9701}{rm}{48}%
}\hskip.1em\allowbreak{} ，将其唯一分解为

\[
\RB{140.840}{7.953}{2.604}{%
\RG{0.000}{0.002}{10.9091}{mi}{118}\RG{8.585}{0.002}{10.9091}{rm}{61}\RG{19.924}{0.002}{10.9091}{mi}{118}\RG{31.815}{0.002}{10.9091}{rm}{43}\RG{42.546}{0.002}{10.9091}{sy}{1}\RG{47.390}{0.002}{10.9091}{sy}{1}\RG{52.245}{0.002}{10.9091}{sy}{1}\RG{57.698}{0.002}{10.9091}{rm}{43}\RG{68.431}{0.002}{10.9091}{mi}{118}\RG{80.857}{0.002}{10.9091}{mi}{59}\RG{107.454}{0.002}{10.9091}{mi}{118}\RG{118.729}{0.002}{10.9091}{sy}{50}\RG{129.032}{0.002}{10.9091}{mi}{86}\RG{138.333}{0.002}{10.9091}{mi}{58}\RG{25.088}{-1.635}{7.9701}{rm}{49}\RG{73.595}{-1.635}{7.9701}{mi}{109}\RG{112.625}{-1.635}{7.9701}{mi}{105}\RG{135.259}{-1.635}{7.9701}{mi}{105}%
}
\]
于是

\[
\RB{166.416}{15.482}{11.387}{%
\RG{87.497}{11.787}{7.9701}{mi}{109}\RG{97.444}{10.519}{9.0569}{ex}{16}\RG{158.500}{10.519}{9.0569}{ex}{17}\RG{84.361}{9.070}{9.0569}{ex}{88}\RG{0.000}{0.003}{10.9091}{rm}{48}\RG{8.367}{0.003}{10.9091}{rm}{61}\RG{19.710}{0.003}{10.9091}{mi}{39}\RG{26.707}{0.003}{10.9091}{rm}{40}\RG{30.861}{0.003}{10.9091}{mi}{118}\RG{36.416}{0.003}{10.9091}{rm}{41}\RG{42.995}{0.003}{10.9091}{sy}{0}\RG{53.904}{0.003}{10.9091}{mi}{21}\RG{64.437}{0.003}{10.9091}{mi}{118}\RG{73.023}{0.003}{10.9091}{rm}{61}\RG{102.853}{0.003}{10.9091}{mi}{39}\RG{109.850}{0.003}{10.9091}{rm}{40}\RG{114.004}{0.003}{10.9091}{mi}{118}\RG{122.242}{0.003}{10.9091}{rm}{41}\RG{128.820}{0.003}{10.9091}{sy}{0}\RG{139.729}{0.003}{10.9091}{mi}{21}\RG{150.262}{0.003}{10.9091}{mi}{118}\RG{163.909}{0.003}{10.9091}{mi}{58}\RG{60.135}{-1.634}{7.9701}{rm}{48}\RG{119.168}{-1.634}{7.9701}{mi}{105}\RG{145.960}{-1.634}{7.9701}{rm}{48}\RG{155.426}{-1.634}{7.9701}{mi}{105}\RG{84.674}{-10.815}{7.9701}{mi}{105}\RG{87.295}{-10.815}{7.9701}{rm}{61}\RG{93.279}{-10.815}{7.9701}{rm}{49}%
}
\]
由\allowbreak{}\hskip.1em \RB{9.024}{7.951}{2.209}{%
\RG{0.000}{-0.000}{10.9091}{mi}{86}\RG{6.227}{-1.636}{7.9701}{mi}{105}%
}\hskip.1em\allowbreak{} 的不变性，每个括号中的向量都属于\allowbreak{}\hskip.1em \RB{9.024}{7.951}{2.209}{%
\RG{0.000}{-0.000}{10.9091}{mi}{86}\RG{6.227}{-1.636}{7.9701}{mi}{105}%
}\hskip.1em\allowbreak{} 。根据直和分解的唯一性，

\[
\RB{155.529}{8.685}{3.235}{%
\RG{0.000}{0.004}{10.9091}{mi}{39}\RG{6.997}{0.004}{10.9091}{rm}{40}\RG{11.151}{0.004}{10.9091}{mi}{118}\RG{19.389}{0.004}{10.9091}{rm}{41}\RG{25.967}{0.004}{10.9091}{sy}{0}\RG{36.876}{0.004}{10.9091}{mi}{21}\RG{47.409}{0.004}{10.9091}{mi}{118}\RG{58.677}{0.004}{10.9091}{rm}{61}\RG{70.012}{0.004}{10.9091}{rm}{48}\RG{97.175}{0.004}{10.9091}{rm}{40}\RG{101.332}{0.004}{10.9091}{rm}{49}\RG{109.700}{0.004}{10.9091}{sy}{20}\RG{121.215}{0.004}{10.9091}{mi}{105}\RG{127.890}{0.004}{10.9091}{sy}{20}\RG{139.405}{0.004}{10.9091}{mi}{109}\RG{148.748}{0.004}{10.9091}{rm}{41}\RG{152.902}{0.004}{10.9091}{mi}{59}\RG{16.315}{-1.633}{7.9701}{mi}{105}\RG{43.107}{-1.633}{7.9701}{rm}{48}\RG{52.574}{-1.633}{7.9701}{mi}{105}%
}
\]
故\allowbreak{}\hskip.1em \RB{36.601}{7.949}{3.537}{%
\RG{0.000}{-0.002}{10.9091}{mi}{118}\RG{5.165}{-1.638}{7.9701}{mi}{105}\RG{11.268}{-0.002}{10.9091}{sy}{50}\RG{21.572}{-0.002}{10.9091}{mi}{86}\RG{27.799}{-1.638}{7.9701}{mi}{105}\RG{30.422}{-1.638}{7.9701}{mi}{59}\RG{32.558}{-1.638}{7.9701}{rm}{48}%
}\hskip.1em\allowbreak{} 。因此反向包含关系也成立，所证等式成立。（2）将\allowbreak{}\hskip.1em \RB{12.813}{8.183}{0.705}{%
\RG{0.000}{0.003}{10.9091}{bb}{67}\RG{7.879}{3.961}{7.9701}{mi}{110}%
}\hskip.1em\allowbreak{} 中的列向量按阶数\allowbreak{}\hskip.1em \RB{48.053}{5.314}{2.603}{%
\RG{0.000}{0.003}{10.9091}{mi}{110}\RG{6.375}{-1.634}{7.9701}{rm}{49}\RG{10.678}{0.003}{10.9091}{mi}{59}\RG{15.456}{0.003}{10.9091}{mi}{58}\RG{20.245}{0.003}{10.9091}{mi}{58}\RG{25.024}{0.003}{10.9091}{mi}{58}\RG{29.813}{0.003}{10.9091}{mi}{59}\RG{34.602}{0.003}{10.9091}{mi}{110}\RG{40.981}{-1.634}{7.9701}{mi}{109}%
}\hskip.1em\allowbreak{} 分块，记\allowbreak{}\hskip.1em \RB{70.551}{8.183}{2.206}{%
\RG{0.000}{0.003}{10.9091}{mi}{19}\RG{3.775}{-1.634}{7.9701}{mi}{105}\RG{10.466}{0.003}{10.9091}{rm}{58}\RG{17.051}{0.003}{10.9091}{bb}{67}\RG{24.929}{3.961}{7.9701}{mi}{110}\RG{29.600}{2.856}{5.9776}{mi}{105}\RG{36.545}{0.003}{10.9091}{sy}{0}\RG{43.221}{0.003}{10.9091}{sy}{33}\RG{57.737}{0.003}{10.9091}{bb}{67}\RG{65.616}{3.961}{7.9701}{mi}{110}%
}\hskip.1em\allowbreak{} 为补零嵌入：\allowbreak{}\hskip.1em \RB{20.659}{8.684}{3.236}{%
\RG{0.000}{0.003}{10.9091}{mi}{19}\RG{3.775}{-1.634}{7.9701}{mi}{105}\RG{6.849}{0.003}{10.9091}{rm}{40}\RG{11.004}{0.003}{10.9091}{mi}{117}\RG{17.083}{0.003}{10.9091}{rm}{41}%
}\hskip.1em\allowbreak{}的第\allowbreak{}\hskip.1em \RB{3.587}{7.670}{0.607}{%
\RG{0.000}{0.003}{10.9091}{mi}{105}%
}\hskip.1em\allowbreak{} 个列块是\allowbreak{}\hskip.1em \RB{6.260}{5.313}{0.607}{%
\RG{0.000}{0.003}{10.9091}{mi}{117}%
}\hskip.1em\allowbreak{} ，其余各列块均为零。由分块矩阵乘法，对任意\allowbreak{}\hskip.1em \RB{34.549}{8.183}{0.934}{%
\RG{0.000}{0.003}{10.9091}{mi}{117}\RG{9.110}{0.003}{10.9091}{sy}{50}\RG{19.413}{0.003}{10.9091}{bb}{67}\RG{27.292}{3.961}{7.9701}{mi}{110}\RG{31.963}{2.856}{5.9776}{mi}{105}%
}\hskip.1em\allowbreak{} ，有

\[
\RB{78.438}{8.679}{3.241}{%
\RG{0.000}{-0.003}{10.9091}{mi}{65}\RG{8.007}{-0.003}{10.9091}{mi}{19}\RG{14.856}{-0.003}{10.9091}{rm}{40}\RG{19.011}{-0.003}{10.9091}{mi}{117}\RG{25.090}{-0.003}{10.9091}{rm}{41}\RG{32.268}{-0.003}{10.9091}{rm}{61}\RG{43.613}{-0.003}{10.9091}{mi}{19}\RG{50.462}{-0.003}{10.9091}{rm}{40}\RG{54.615}{-0.003}{10.9091}{mi}{65}\RG{65.697}{-0.003}{10.9091}{mi}{117}\RG{71.776}{-0.003}{10.9091}{rm}{41}\RG{75.931}{-0.003}{10.9091}{mi}{58}\RG{11.783}{-1.639}{7.9701}{mi}{105}\RG{47.388}{-1.639}{7.9701}{mi}{105}\RG{62.624}{-1.639}{7.9701}{mi}{105}%
}
\]
取\allowbreak{}\hskip.1em \RB{48.685}{8.683}{3.237}{%
\RG{0.000}{0.001}{10.9091}{mi}{11}\RG{1.213}{0.798}{9.0569}{ex}{101}\RG{6.795}{-1.636}{7.9701}{mi}{105}\RG{12.898}{0.001}{10.9091}{rm}{61}\RG{24.237}{0.001}{10.9091}{mi}{19}\RG{28.012}{-1.636}{7.9701}{mi}{105}\RG{31.086}{0.001}{10.9091}{rm}{40}\RG{35.240}{0.001}{10.9091}{mi}{11}\RG{42.034}{-1.636}{7.9701}{mi}{105}\RG{45.108}{0.001}{10.9091}{rm}{41}%
}\hskip.1em\allowbreak{} ，则

\[
\RB{153.098}{8.683}{3.237}{%
\RG{9.220}{0.798}{9.0569}{ex}{101}\RG{141.935}{0.798}{9.0569}{ex}{101}\RG{0.000}{0.001}{10.9091}{mi}{65}\RG{8.007}{0.001}{10.9091}{mi}{11}\RG{20.905}{0.001}{10.9091}{rm}{61}\RG{32.245}{0.001}{10.9091}{mi}{19}\RG{39.094}{0.001}{10.9091}{rm}{40}\RG{43.247}{0.001}{10.9091}{mi}{65}\RG{54.328}{0.001}{10.9091}{mi}{11}\RG{64.196}{0.001}{10.9091}{rm}{41}\RG{71.375}{0.001}{10.9091}{rm}{61}\RG{82.719}{0.001}{10.9091}{mi}{19}\RG{89.568}{0.001}{10.9091}{rm}{40}\RG{93.723}{0.001}{10.9091}{mi}{21}\RG{103.027}{0.001}{10.9091}{mi}{11}\RG{112.895}{0.001}{10.9091}{rm}{41}\RG{120.074}{0.001}{10.9091}{rm}{61}\RG{131.418}{0.001}{10.9091}{mi}{21}\RG{140.723}{0.001}{10.9091}{mi}{11}\RG{150.591}{0.001}{10.9091}{mi}{58}\RG{14.802}{-1.636}{7.9701}{mi}{105}\RG{36.019}{-1.636}{7.9701}{mi}{105}\RG{51.255}{-1.636}{7.9701}{mi}{105}\RG{61.123}{-1.636}{7.9701}{mi}{105}\RG{86.495}{-1.636}{7.9701}{mi}{105}\RG{99.954}{-1.636}{7.9701}{mi}{105}\RG{109.821}{-1.636}{7.9701}{mi}{105}\RG{137.649}{-1.636}{7.9701}{mi}{105}\RG{147.517}{-1.636}{7.9701}{mi}{105}%
}
\]
又因\allowbreak{}\hskip.1em \RB{29.620}{8.313}{2.854}{%
\RG{0.000}{0.003}{10.9091}{mi}{11}\RG{6.795}{-1.634}{7.9701}{mi}{105}\RG{12.898}{0.003}{10.9091}{sy}{54}\RG{12.898}{0.003}{10.9091}{rm}{61}\RG{24.233}{0.003}{10.9091}{rm}{48}%
}\hskip.1em\allowbreak{} ，故\allowbreak{}\hskip.1em \RB{29.620}{8.313}{2.854}{%
\RG{0.000}{0.003}{10.9091}{mi}{11}\RG{1.213}{0.800}{9.0569}{ex}{101}\RG{6.794}{-1.634}{7.9701}{mi}{105}\RG{12.898}{0.003}{10.9091}{sy}{54}\RG{12.898}{0.003}{10.9091}{rm}{61}\RG{24.233}{0.003}{10.9091}{rm}{48}%
}\hskip.1em\allowbreak{} ，它确为所需特征向量。

% Original page 7.
（3）记

\[
\RB{211.842}{8.685}{3.293}{%
\RG{0.000}{0.003}{10.9091}{mi}{69}\RG{14.011}{0.003}{10.9091}{rm}{61}\RG{25.345}{0.003}{10.9091}{rm}{107}\RG{30.702}{0.003}{10.9091}{rm}{101}\RG{35.447}{0.003}{10.9091}{rm}{114}\RG{39.604}{0.003}{10.9091}{rm}{40}\RG{43.745}{0.003}{10.9091}{mi}{65}\RG{57.251}{0.003}{10.9091}{sy}{0}\RG{68.160}{0.003}{10.9091}{mi}{21}\RG{78.694}{0.003}{10.9091}{mi}{73}\RG{91.403}{0.003}{10.9091}{rm}{41}\RG{95.556}{0.003}{10.9091}{mi}{59}\RG{122.153}{0.003}{10.9091}{mi}{69}\RG{133.737}{0.003}{10.9091}{rm}{61}\RG{145.072}{0.003}{10.9091}{rm}{107}\RG{150.428}{0.003}{10.9091}{rm}{101}\RG{155.174}{0.003}{10.9091}{rm}{114}\RG{159.330}{0.003}{10.9091}{rm}{40}\RG{163.472}{0.003}{10.9091}{mi}{65}\RG{173.904}{0.003}{10.9091}{sy}{0}\RG{184.813}{0.003}{10.9091}{mi}{21}\RG{195.345}{0.003}{10.9091}{mi}{73}\RG{205.181}{0.003}{10.9091}{rm}{41}\RG{209.335}{0.003}{10.9091}{mi}{58}\RG{7.907}{-1.634}{7.9701}{mi}{105}\RG{51.753}{-1.634}{7.9701}{mi}{105}\RG{74.391}{-1.634}{7.9701}{rm}{48}\RG{83.405}{-1.634}{7.9701}{mi}{110}\RG{191.044}{-1.634}{7.9701}{rm}{48}\RG{200.056}{-1.634}{7.9701}{mi}{110}\RG{88.075}{-2.738}{5.9776}{mi}{105}%
}
\]
对分块列向量\allowbreak{}\hskip.1em \RB{83.264}{9.434}{3.272}{%
\RG{0.000}{-0.003}{10.9091}{mi}{120}\RG{9.101}{-0.003}{10.9091}{rm}{61}\RG{20.435}{-0.003}{10.9091}{rm}{40}\RG{24.593}{-0.003}{10.9091}{mi}{120}\RG{30.663}{3.956}{7.9701}{rm}{84}\RG{30.663}{-2.700}{7.9701}{rm}{49}\RG{36.374}{-0.003}{10.9091}{mi}{59}\RG{41.152}{-0.003}{10.9091}{mi}{58}\RG{45.941}{-0.003}{10.9091}{mi}{58}\RG{50.730}{-0.003}{10.9091}{mi}{58}\RG{55.508}{-0.003}{10.9091}{mi}{59}\RG{60.297}{-0.003}{10.9091}{mi}{120}\RG{66.370}{3.956}{7.9701}{rm}{84}\RG{66.370}{-2.700}{7.9701}{mi}{109}\RG{73.633}{-0.003}{10.9091}{rm}{41}\RG{77.786}{3.956}{7.9701}{rm}{84}%
}\hskip.1em\allowbreak{} ，有

\[
\RB{377.764}{8.684}{3.293}{%
\RG{0.000}{0.002}{10.9091}{mi}{120}\RG{9.100}{0.002}{10.9091}{sy}{50}\RG{19.404}{0.002}{10.9091}{mi}{69}\RG{41.889}{0.002}{10.9091}{sy}{40}\RG{50.987}{0.002}{10.9091}{sy}{41}\RG{75.829}{0.002}{10.9091}{rm}{40}\RG{79.983}{0.002}{10.9091}{mi}{65}\RG{93.488}{0.002}{10.9091}{sy}{0}\RG{104.397}{0.002}{10.9091}{mi}{21}\RG{114.930}{0.002}{10.9091}{mi}{73}\RG{127.640}{0.002}{10.9091}{rm}{41}\RG{131.794}{0.002}{10.9091}{mi}{120}\RG{143.968}{0.002}{10.9091}{rm}{61}\RG{155.303}{0.002}{10.9091}{rm}{48}\RG{171.557}{0.002}{10.9091}{rm}{40}\RG{175.714}{0.002}{10.9091}{rm}{49}\RG{184.082}{0.002}{10.9091}{sy}{20}\RG{195.596}{0.002}{10.9091}{mi}{105}\RG{202.271}{0.002}{10.9091}{sy}{20}\RG{213.786}{0.002}{10.9091}{mi}{109}\RG{223.129}{0.002}{10.9091}{rm}{41}\RG{241.223}{0.002}{10.9091}{sy}{40}\RG{250.321}{0.002}{10.9091}{sy}{41}\RG{275.162}{0.002}{10.9091}{mi}{120}\RG{287.336}{0.002}{10.9091}{sy}{50}\RG{297.639}{0.002}{10.9091}{mi}{69}\RG{319.529}{0.002}{10.9091}{rm}{40}\RG{323.685}{0.002}{10.9091}{rm}{49}\RG{332.055}{0.002}{10.9091}{sy}{20}\RG{343.570}{0.002}{10.9091}{mi}{105}\RG{350.244}{0.002}{10.9091}{sy}{20}\RG{361.759}{0.002}{10.9091}{mi}{109}\RG{371.103}{0.002}{10.9091}{rm}{41}\RG{375.256}{0.002}{10.9091}{mi}{58}\RG{87.990}{-1.635}{7.9701}{mi}{105}\RG{110.628}{-1.635}{7.9701}{rm}{48}\RG{119.642}{-1.635}{7.9701}{mi}{110}\RG{137.864}{-1.635}{7.9701}{mi}{105}\RG{281.232}{-1.635}{7.9701}{mi}{105}\RG{305.546}{-1.635}{7.9701}{mi}{105}\RG{124.313}{-2.739}{5.9776}{mi}{105}%
}
\]
若\allowbreak{}\hskip.1em \RB{67.247}{8.679}{3.230}{%
\RG{0.000}{-0.003}{10.9091}{mi}{105}\RG{6.675}{-0.003}{10.9091}{sy}{50}\RG{7.942}{-0.003}{10.9091}{mi}{61}\RG{16.973}{-0.003}{10.9091}{sy}{102}\RG{22.432}{-0.003}{10.9091}{mi}{105}\RG{26.076}{-1.639}{7.9701}{rm}{49}\RG{30.378}{-0.003}{10.9091}{mi}{59}\RG{35.156}{-0.003}{10.9091}{mi}{58}\RG{39.945}{-0.003}{10.9091}{mi}{58}\RG{44.724}{-0.003}{10.9091}{mi}{58}\RG{49.513}{-0.003}{10.9091}{mi}{59}\RG{54.302}{-0.003}{10.9091}{mi}{105}\RG{57.948}{-1.639}{7.9701}{mi}{114}\RG{62.089}{-0.003}{10.9091}{sy}{103}%
}\hskip.1em\allowbreak{} ，则\allowbreak{}\hskip.1em \RB{10.274}{8.068}{2.291}{%
\RG{0.000}{-0.003}{10.9091}{mi}{21}\RG{6.231}{-1.639}{7.9701}{rm}{48}%
}\hskip.1em\allowbreak{} 不是\allowbreak{}\hskip.1em \RB{10.804}{8.286}{2.212}{%
\RG{0.000}{-0.003}{10.9091}{mi}{65}\RG{8.007}{-1.639}{7.9701}{mi}{105}%
}\hskip.1em\allowbreak{} 的特征值，故\allowbreak{}\hskip.1em \RB{41.304}{8.679}{3.230}{%
\RG{0.000}{-0.003}{10.9091}{mi}{69}\RG{7.907}{-1.639}{7.9701}{mi}{105}\RG{14.011}{-0.003}{10.9091}{rm}{61}\RG{25.350}{-0.003}{10.9091}{sy}{102}\RG{30.805}{-0.003}{10.9091}{rm}{48}\RG{36.145}{-0.003}{10.9091}{sy}{103}%
}\hskip.1em\allowbreak{} 。对每个\allowbreak{}\hskip.1em \RB{45.501}{8.068}{1.997}{%
\RG{0.000}{-0.003}{10.9091}{rm}{49}\RG{8.372}{-0.003}{10.9091}{sy}{20}\RG{19.886}{-0.003}{10.9091}{mi}{107}\RG{28.839}{-0.003}{10.9091}{sy}{20}\RG{40.354}{-0.003}{10.9091}{mi}{114}%
}\hskip.1em\allowbreak{} ，取\allowbreak{}\hskip.1em \RB{14.300}{7.926}{3.472}{%
\RG{0.000}{-0.003}{10.9091}{mi}{69}\RG{7.907}{-1.639}{7.9701}{mi}{105}\RG{10.528}{-2.917}{5.9776}{mi}{107}%
}\hskip.1em\allowbreak{} 的一组基

\[
\RB{151.685}{8.068}{3.537}{%
\RG{0.000}{-0.003}{10.9091}{mi}{11}\RG{17.436}{-0.003}{10.9091}{mi}{59}\RG{22.215}{-0.003}{10.9091}{mi}{58}\RG{27.004}{-0.003}{10.9091}{mi}{58}\RG{31.793}{-0.003}{10.9091}{mi}{58}\RG{36.571}{-0.003}{10.9091}{mi}{59}\RG{41.360}{-0.003}{10.9091}{mi}{11}\RG{63.068}{-0.003}{10.9091}{mi}{59}\RG{89.665}{-0.003}{10.9091}{mi}{100}\RG{102.907}{-0.003}{10.9091}{rm}{61}\RG{114.242}{-0.003}{10.9091}{rm}{100}\RG{120.176}{-0.003}{10.9091}{rm}{105}\RG{123.144}{-0.003}{10.9091}{rm}{109}\RG{133.919}{-0.003}{10.9091}{mi}{69}\RG{149.057}{-0.003}{10.9091}{mi}{59}\RG{6.795}{-1.639}{7.9701}{mi}{107}\RG{10.990}{-1.639}{7.9701}{mi}{59}\RG{13.135}{-1.639}{7.9701}{rm}{49}\RG{48.157}{-1.639}{7.9701}{mi}{107}\RG{52.352}{-1.639}{7.9701}{mi}{59}\RG{54.490}{-1.639}{7.9701}{mi}{100}\RG{95.223}{-1.639}{7.9701}{mi}{107}\RG{141.826}{-1.639}{7.9701}{mi}{105}\RG{58.458}{-2.917}{5.9776}{mi}{107}\RG{144.447}{-2.917}{5.9776}{mi}{107}%
}
\]
并按（2）延拓为\allowbreak{}\hskip.1em \RB{62.230}{7.837}{3.534}{%
\RG{0.000}{0.001}{10.9091}{mi}{11}\RG{1.213}{0.798}{9.0569}{ex}{101}\RG{6.794}{-1.635}{7.9701}{mi}{107}\RG{10.989}{-1.635}{7.9701}{mi}{59}\RG{13.134}{-1.635}{7.9701}{rm}{49}\RG{17.436}{0.001}{10.9091}{mi}{59}\RG{22.215}{0.001}{10.9091}{mi}{58}\RG{27.004}{0.001}{10.9091}{mi}{58}\RG{31.782}{0.001}{10.9091}{mi}{58}\RG{36.571}{0.001}{10.9091}{mi}{59}\RG{41.363}{0.001}{10.9091}{mi}{11}\RG{42.575}{0.798}{9.0569}{ex}{101}\RG{48.156}{-1.635}{7.9701}{mi}{107}\RG{52.352}{-1.635}{7.9701}{mi}{59}\RG{54.489}{-1.635}{7.9701}{mi}{100}\RG{58.458}{-2.914}{5.9776}{mi}{107}%
}\hskip.1em\allowbreak{} 。由（2），这些延拓向量均属于\allowbreak{}\hskip.1em \RB{8.715}{7.930}{0.500}{%
\RG{0.000}{0.001}{10.9091}{mi}{69}%
}\hskip.1em\allowbreak{} 。任取\allowbreak{}\hskip.1em \RB{28.118}{7.930}{0.935}{%
\RG{0.000}{0.001}{10.9091}{mi}{120}\RG{9.101}{0.001}{10.9091}{sy}{50}\RG{19.404}{0.001}{10.9091}{mi}{69}%
}\hskip.1em\allowbreak{} ，其余块皆为零，而每个\allowbreak{}\hskip.1em \RB{12.462}{5.312}{3.468}{%
\RG{0.000}{0.001}{10.9091}{mi}{120}\RG{6.070}{-1.635}{7.9701}{mi}{105}\RG{8.691}{-2.914}{5.9776}{mi}{107}%
}\hskip.1em\allowbreak{} 均可用\allowbreak{}\hskip.1em \RB{62.229}{5.312}{3.534}{%
\RG{0.000}{0.001}{10.9091}{mi}{11}\RG{6.794}{-1.636}{7.9701}{mi}{107}\RG{10.989}{-1.636}{7.9701}{mi}{59}\RG{13.134}{-1.636}{7.9701}{rm}{49}\RG{17.435}{0.001}{10.9091}{mi}{59}\RG{22.214}{0.001}{10.9091}{mi}{58}\RG{27.003}{0.001}{10.9091}{mi}{58}\RG{31.792}{0.001}{10.9091}{mi}{58}\RG{36.570}{0.001}{10.9091}{mi}{59}\RG{41.359}{0.001}{10.9091}{mi}{11}\RG{48.156}{-1.636}{7.9701}{mi}{107}\RG{52.352}{-1.636}{7.9701}{mi}{59}\RG{54.489}{-1.636}{7.9701}{mi}{100}\RG{58.457}{-2.914}{5.9776}{mi}{107}%
}\hskip.1em\allowbreak{} 线性表示，所以这些延拓向量张成\allowbreak{}\hskip.1em \RB{8.715}{7.930}{0.500}{%
\RG{0.000}{0.001}{10.9091}{mi}{69}%
}\hskip.1em\allowbreak{} 。

若

\[
\RB{84.360}{19.418}{10.014}{%
\RG{18.338}{13.890}{7.9701}{mi}{100}\RG{5.175}{13.416}{7.9701}{mi}{114}\RG{22.299}{12.612}{5.9776}{mi}{107}\RG{0.478}{10.699}{9.0569}{ex}{88}\RG{15.856}{10.699}{9.0569}{ex}{88}\RG{46.504}{2.430}{9.0569}{ex}{101}\RG{30.756}{1.633}{10.9091}{mi}{99}\RG{45.291}{1.633}{10.9091}{mi}{11}\RG{65.054}{1.633}{10.9091}{rm}{61}\RG{76.388}{1.633}{10.9091}{rm}{48}\RG{81.733}{1.633}{10.9091}{mi}{59}\RG{35.354}{-0.004}{7.9701}{mi}{107}\RG{39.549}{-0.004}{7.9701}{mi}{59}\RG{41.686}{-0.004}{7.9701}{mi}{96}\RG{52.085}{-0.004}{7.9701}{mi}{107}\RG{56.281}{-0.004}{7.9701}{mi}{59}\RG{58.418}{-0.004}{7.9701}{mi}{96}\RG{0.000}{-9.427}{7.9701}{mi}{107}\RG{4.201}{-9.427}{7.9701}{rm}{61}\RG{10.186}{-9.427}{7.9701}{rm}{49}\RG{15.906}{-9.427}{7.9701}{mi}{96}\RG{19.051}{-9.427}{7.9701}{rm}{61}\RG{25.036}{-9.427}{7.9701}{rm}{49}%
}
\]
比较第\allowbreak{}\hskip.1em \RB{8.006}{7.665}{2.211}{%
\RG{0.000}{-0.002}{10.9091}{mi}{105}\RG{3.645}{-1.638}{7.9701}{mi}{107}%
}\hskip.1em\allowbreak{} 块，得到

\[
\RB{68.384}{19.424}{10.009}{%
\RG{2.482}{13.895}{7.9701}{mi}{100}\RG{6.443}{12.617}{5.9776}{mi}{107}\RG{0.000}{10.705}{9.0569}{ex}{88}\RG{14.900}{1.638}{10.9091}{mi}{99}\RG{29.435}{1.638}{10.9091}{mi}{11}\RG{49.197}{1.638}{10.9091}{rm}{61}\RG{60.532}{1.638}{10.9091}{rm}{48}\RG{65.876}{1.638}{10.9091}{mi}{58}\RG{19.497}{0.002}{7.9701}{mi}{107}\RG{23.692}{0.002}{7.9701}{mi}{59}\RG{25.830}{0.002}{7.9701}{mi}{96}\RG{36.229}{0.002}{7.9701}{mi}{107}\RG{40.424}{0.002}{7.9701}{mi}{59}\RG{42.562}{0.002}{7.9701}{mi}{96}\RG{0.050}{-9.422}{7.9701}{mi}{96}\RG{3.195}{-9.422}{7.9701}{rm}{61}\RG{9.179}{-9.422}{7.9701}{rm}{49}%
}
\]
由每组基的线性无关性，所有\allowbreak{}\hskip.1em \RB{14.422}{5.315}{3.531}{%
\RG{0.000}{0.004}{10.9091}{mi}{99}\RG{4.597}{-1.633}{7.9701}{mi}{107}\RG{8.792}{-1.633}{7.9701}{mi}{59}\RG{10.930}{-1.633}{7.9701}{mi}{96}%
}\hskip.1em\allowbreak{} 都为零。因此延拓后的并集线性无关，确为\allowbreak{}\hskip.1em \RB{8.715}{7.933}{0.500}{%
\RG{0.000}{0.004}{10.9091}{mi}{69}%
}\hskip.1em\allowbreak{} 的一组基。

小结：不变直和分解会诱导同一特征值的特征子空间的直和分解；分块对角矩阵正是这一事实的坐标形式。尤其有

\[
\RB{205.711}{15.478}{11.392}{%
\RG{101.438}{11.783}{7.9701}{mi}{109}\RG{98.302}{9.065}{9.0569}{ex}{88}\RG{0.000}{-0.002}{10.9091}{rm}{100}\RG{5.935}{-0.002}{10.9091}{rm}{105}\RG{8.902}{-0.002}{10.9091}{rm}{109}\RG{19.669}{-0.002}{10.9091}{rm}{107}\RG{25.025}{-0.002}{10.9091}{rm}{101}\RG{29.771}{-0.002}{10.9091}{rm}{114}\RG{33.927}{-0.002}{10.9091}{rm}{40}\RG{38.070}{-0.002}{10.9091}{mi}{65}\RG{48.502}{-0.002}{10.9091}{sy}{0}\RG{59.411}{-0.002}{10.9091}{mi}{21}\RG{69.944}{-0.002}{10.9091}{mi}{73}\RG{79.779}{-0.002}{10.9091}{rm}{41}\RG{86.957}{-0.002}{10.9091}{rm}{61}\RG{113.202}{-0.002}{10.9091}{rm}{100}\RG{119.136}{-0.002}{10.9091}{rm}{105}\RG{122.104}{-0.002}{10.9091}{rm}{109}\RG{132.871}{-0.002}{10.9091}{rm}{107}\RG{138.227}{-0.002}{10.9091}{rm}{101}\RG{142.973}{-0.002}{10.9091}{rm}{114}\RG{147.129}{-0.002}{10.9091}{rm}{40}\RG{151.272}{-0.002}{10.9091}{mi}{65}\RG{164.778}{-0.002}{10.9091}{sy}{0}\RG{175.687}{-0.002}{10.9091}{mi}{21}\RG{186.220}{-0.002}{10.9091}{mi}{73}\RG{198.929}{-0.002}{10.9091}{rm}{41}\RG{203.084}{-0.002}{10.9091}{mi}{59}\RG{65.642}{-1.638}{7.9701}{rm}{48}\RG{74.655}{-1.638}{7.9701}{mi}{110}\RG{159.280}{-1.638}{7.9701}{mi}{105}\RG{181.918}{-1.638}{7.9701}{rm}{48}\RG{190.931}{-1.638}{7.9701}{mi}{110}\RG{195.602}{-2.743}{5.9776}{mi}{105}\RG{98.614}{-10.820}{7.9701}{mi}{105}\RG{101.235}{-10.820}{7.9701}{rm}{61}\RG{107.220}{-10.820}{7.9701}{rm}{49}%
}
\]
即该特征值的几何重数等于各块中对应几何重数之和。


\end{document}
